\chapter{SİSTEM MİMARİSİ, MODEL VE METODOLOJİ}
\label{chap:sistem}
% T-tipi faz kolu için ortak vektör çizim tabanı.
% Argümanlar: Q1, Q2, orta nokta tarafı nötr anahtar (Q4, düğüm q3), faz tarafı nötr anahtar (Q3, düğüm q4) çizgi renkleri ve ek TikZ yolu.
% Adlandırma (09.11.2026, K-31): tümleyen çiftler Q1--Q3 ve Q2--Q4; düğüm adları konumu gösterir.
\newcommand{\ttypevector}[5]{%
\begin{tikzpicture}[
  x=0.76cm,y=0.56cm,
  terminal/.style={circle,fill=black,inner sep=1.25pt},
  sw/.style={rounded corners=1pt,minimum width=7mm,minimum height=8mm,fill=white,font=\scriptsize},
  bus/.style={line width=0.8pt},
  current/.style={-{Latex[length=2mm]},red!80!black,line width=1.7pt},
  activepath/.style={blue!70!black,line width=2.0pt}
]
  % Bölünmüş DA bara ve kapasitörler
  \draw[bus] (0,5) -- (4,5);
  \draw[bus] (0,-1) -- (4,-1);
  \draw[bus] (0,5) -- (0,3.65);
  \draw[bus] (-0.32,3.65) -- (0.32,3.65);
  \draw[bus] (-0.32,3.35) -- (0.32,3.35);
  \draw[bus] (0,3.35) -- (0,2);
  \draw[bus] (0,2) -- (0,0.65);
  \draw[bus] (-0.32,0.65) -- (0.32,0.65);
  \draw[bus] (-0.32,0.35) -- (0.32,0.35);
  \draw[bus] (0,0.35) -- (0,-1);
  \node[font=\scriptsize,left] at (0,5) {$+V_{\mathrm{dc}}/2$};
  \node[font=\scriptsize,left] at (0,2) {$N$};
  \node[font=\scriptsize,left] at (0,-1) {$-V_{\mathrm{dc}}/2$};
  \node[font=\scriptsize,left] at (-0.38,3.5) {$C_A$};
  \node[font=\scriptsize,left] at (-0.38,0.5) {$C_B$};

  % Ana anahtarlar ve çift yönlü nötr kol
  \node[sw,draw=#1,line width=1.25pt] (q1) at (4,3.65) {$Q_1$};
  \node[sw,draw=#2,line width=1.25pt] (q2) at (4,0.35) {$Q_2$};
  \node[sw,draw=#3,line width=1.25pt] (q3) at (1.55,2) {$Q_4$};
  \node[sw,draw=#4,line width=1.25pt] (q4) at (2.75,2) {$Q_3$};
  \draw[bus] (4,5) -- (q1.north);
  \draw[bus] (q1.south) -- (4,2);
  \draw[bus] (4,2) -- (q2.north);
  \draw[bus] (q2.south) -- (4,-1);
  \draw[bus] (0,2) -- (q3.west);
  \draw[bus] (q3.east) -- (q4.west);
  \draw[bus] (q4.east) -- (5.15,2);
  \node[terminal] at (0,2) {};
  \node[terminal] (out) at (4,2) {};
  \draw[-{Latex[length=1.8mm]},thick] (out) -- (5.15,2);
  \node[font=\scriptsize,right] at (5.15,2) {$v_{xN},\ i_x$};
  #5
\end{tikzpicture}%
}

% Faz başına tam T-tipi güç devresi: bölünmüş DA bara, dört IGBT ve LCL filtre.
% Filtre kapasitörü modelle tutarlı olarak C_f adıyla gösterilir.
\newcommand{\ttypefulltopology}{%
\begin{tikzpicture}[
  x=0.88cm,y=0.68cm,
  wire/.style={line width=0.9pt},
  coilpath/.style={decorate,decoration={coil,aspect=0.48,segment length=2.8mm,amplitude=1.8mm}},
  igbt/.style={nigbt,bodydiode,anchor=C,tr circle=false},
  terminal/.style={circle,fill=black,inner sep=1.35pt},
  label/.style={font=\scriptsize,align=center}
]
  \ctikzset{transistors/scale=0.85,transistor bodydiode/scale=0.40}
  % Bölünmüş DA bara ve orta nokta
  \draw[wire] (0,5) -- (5.3,5);
  \draw[wire] (0,-1) -- (5.3,-1);
  \draw[wire] (0,5) -- (0,3.65);
  \draw[wire] (-0.38,3.65) -- (0.38,3.65);
  \draw[wire] (-0.38,3.35) -- (0.38,3.35);
  \draw[wire] (0,3.35) -- (0,2);
  \draw[wire] (0,2) -- (0,0.65);
  \draw[wire] (-0.38,0.65) -- (0.38,0.65);
  \draw[wire] (-0.38,0.35) -- (0.38,0.35);
  \draw[wire] (0,0.35) -- (0,-1);
  \node[label,left] at (0,5) {$\mathrm{DC}^{+}$};
  \node[label,above left=0.7mm and 0.8mm] at (0,2) {$N$};
  \node[label,left] at (0,-1) {$\mathrm{DC}^{-}$};
  \node[label,left=1mm] at (-0.38,3.50) {$C_A$};
  \node[label,left=1mm] at (-0.38,0.50) {$C_B$};
  \node[terminal] at (0,2) {};

  % Dış faz kolu ve emetörden kolektöre yönelen antiparalel diyotlar.
  \node[igbt] (q1) at (5.3,4.43) {};
  \node[igbt] (q2) at (5.3,1.13) {};
  \draw[wire] (5.3,5) -- (q1.C);
  \draw[wire] (q1.E) -- (5.3,2);
  \draw[wire] (5.3,2) -- (q2.C);
  \draw[wire] (q2.E) -- (5.3,-1);
  \node[label,right] at (6.15,3.94) {$Q_1$};
  \node[label,right] at (6.15,3.44) {$DS_1$};
  \node[label,right] at (6.15,0.64) {$Q_2$};
  \node[label,right] at (6.15,0.14) {$DS_2$};
  \node[terminal] (pole) at (5.3,2) {};

  % Nötr kol: Q4 ve Q3'ün emetörleri ortada birleşir.
  \node[igbt,rotate=90] (q4) at (1.02,2) {};
  \node[igbt,rotate=-90,xscale=-1] (q3) at (4.40,2) {};
  \draw[wire] (0,2) -- (q4.C);
  \draw[wire] (q4.E) -- (q3.E);
  \draw[wire] (q3.C) -- (pole);
  \node[label,below] at (1.65,0.98) {$Q_4$};
  \node[label,below] at (3.78,0.98) {$Q_3$};
  \node[label,above] at (1.65,2.77) {$DS_4$};
  \node[label,above] at (3.78,2.77) {$DS_3$};

  % Faz başına LCL filtre ve şebeke/yük terminali
  \coordinate (l1a) at (6.35,2);
  \coordinate (l1b) at (7.45,2);
  \coordinate (cnode) at (8.15,2);
  \coordinate (l2a) at (8.85,2);
  \coordinate (l2b) at (9.95,2);
  \draw[wire] (pole) -- (l1a);
  \draw[coilpath,line width=0.9pt] (l1a) -- (l1b);
  \draw[wire] (l1b) -- (cnode) -- (l2a);
  \draw[coilpath,line width=0.9pt] (l2a) -- (l2b);
  \draw[wire,-{Latex[length=1.8mm]}] (l2b) -- (11.1,2);
  \node[label,above=1.5mm] at ($(l1a)!0.5!(l1b)$) {$L_1$};
  \node[label,above=1.5mm] at ($(l2a)!0.5!(l2b)$) {$L_2$};
  \node[label,right] at (11.1,2) {$v_g,\ i_2$};
  \node[terminal] at (cnode) {};
  \draw[wire] (cnode) -- (8.15,1.25);
  \draw[wire] (7.78,1.25) -- (8.52,1.25);
  \draw[wire] (7.78,0.95) -- (8.52,0.95);
  \draw[wire] (8.15,0.95) -- (8.15,-1.45) -- (-1.20,-1.45);
  \draw[wire] (-1.20,-1.45) -- (-1.20,2) -- (0,2);
  \node[label,right=1mm] at (8.52,1.10) {$C_f$};
  \draw[-{Latex[length=1.8mm]},red!80!black,line width=1.1pt] (6.0,2.30) -- (7.0,2.30);
  \node[label,above] at (6.5,2.30) {$i_1$};
\end{tikzpicture}%
}

% Senkron-buck eşdeğeri: üst ray, alt ray, üst ve alt anahtar etiketleri,
% enerji aktaran anahtarın konumu (1: üst, 0: alt). Yeşil anahtar ve mavi yol aktif enerji aktarımıdır (K-32).
\newcommand{\syncbuckvector}[5]{%
\begin{tikzpicture}[
  x=0.82cm,y=0.65cm,
  sw/.style={draw,rounded corners=1pt,minimum width=8mm,minimum height=8mm,fill=white,font=\scriptsize},
  coilpath/.style={decorate,decoration={coil,aspect=0.48,segment length=2.7mm,amplitude=1.8mm}},
  current/.style={-{Latex[length=2mm]},red!80!black,line width=1.4pt}
]
  \draw[line width=0.8pt] (0,3.2) -- (1.8,3.2);
  \draw[line width=0.8pt] (0,0) -- (1.8,0);
  \node[font=\scriptsize,left] at (0,3.2) {#1};
  \node[font=\scriptsize,left] at (0,0) {#2};
  \ifnum#5=1
    \node[sw,draw=green!55!black,line width=1.2pt] (qh) at (1.8,2.45) {#3};
    \node[sw,draw=gray!65,line width=1.2pt] (ql) at (1.8,0.75) {#4};
  \else
    \node[sw,draw=gray!65,line width=1.2pt] (qh) at (1.8,2.45) {#3};
    \node[sw,draw=green!55!black,line width=1.2pt] (ql) at (1.8,0.75) {#4};
  \fi
  \draw[line width=0.8pt] (1.8,3.2) -- (qh.north);
  \draw[line width=0.8pt] (qh.south) -- (1.8,1.6);
  \draw[line width=0.8pt] (1.8,1.6) -- (ql.north);
  \draw[line width=0.8pt] (ql.south) -- (1.8,0);
  \fill (1.8,1.6) circle (1.5pt);
  \coordinate (l1a) at (2.45,1.6);
  \coordinate (l1b) at (3.35,1.6);
  \coordinate (cnode) at (3.88,1.6);
  \coordinate (l2a) at (4.30,1.6);
  \coordinate (l2b) at (5.20,1.6);
  \draw[line width=0.8pt] (1.8,1.6) -- (l1a);
  \draw[coilpath,line width=0.9pt] (l1a) -- (l1b);
  \draw[line width=0.8pt] (l1b) -- (cnode) -- (l2a);
  \draw[coilpath,line width=0.9pt] (l2a) -- (l2b);
  \ifnum#5=1
    \draw[current] (l2b) -- (6.0,1.6);
  \else
    \draw[current] (6.0,1.6) -- (l2b);
  \fi
  \node[font=\scriptsize,above=2mm] at ($(l1a)!0.5!(l1b)$) {$L_1$};
  \node[font=\scriptsize,above=2mm] at ($(l2a)!0.5!(l2b)$) {$L_2$};
  \node[font=\scriptsize,right] at (6.0,1.6) {$v_g,\ i_L$};
  \fill (cnode) circle (1.3pt);
  \draw[line width=0.8pt] (cnode) -- (3.88,1.0);
  \draw[line width=0.8pt] (3.55,1.0) -- (4.21,1.0);
  \draw[line width=0.8pt] (3.55,0.75) -- (4.21,0.75);
  \draw[line width=0.8pt] (3.88,0.75) -- (3.88,0.2);
  \draw[line width=0.8pt] (3.58,0.2) -- (4.18,0.2);
  \draw[line width=0.8pt] (3.68,0.05) -- (4.08,0.05);
  \node[font=\scriptsize,right] at (4.18,0.88) {$C_f$};
  \ifnum#5=1
    \draw[blue!70!black,line width=1.8pt] (0.5,3.2) -| (qh.north);
    \draw[blue!70!black,line width=1.8pt] (qh.south) -- (1.8,1.6) -- (l1a);
  \else
    \draw[blue!70!black,line width=1.8pt] (l1a) -- (1.8,1.6) -- (ql.north);
    \draw[blue!70!black,line width=1.8pt] (ql.south) -- (1.8,0) -- (0.5,0);
  \fi
\end{tikzpicture}%
}

% Susturmalı yarım periyot eşdeğeri: bir aktif anahtar ve tek yönlü serbest dolaşım yolu.
% Argümanlar: üst ray, alt ray, aktif anahtar, diyot etiketi, anahtar konumu (1: üst, 0: alt).
% Diyot her iki yarım periyotta yukarı yönlüdür: pozitifte N'den kutba, negatifte kutuptan N'ye.
\newcommand{\susturmabuckvector}[5]{%
\begin{tikzpicture}[
  x=0.82cm,y=0.65cm,
  sw/.style={draw,rounded corners=1pt,minimum width=8mm,minimum height=8mm,fill=white,font=\scriptsize},
  coilpath/.style={decorate,decoration={coil,aspect=0.48,segment length=2.7mm,amplitude=1.8mm}},
  current/.style={-{Latex[length=2mm]},red!80!black,line width=1.4pt},
  freewheel/.style={orange!85!black,line width=1.6pt,dash pattern=on 3pt off 1.6pt}
]
  \draw[line width=0.8pt] (0,3.2) -- (1.8,3.2);
  \draw[line width=0.8pt] (0,0) -- (1.8,0);
  \node[font=\scriptsize,left] at (0,3.2) {#1};
  \node[font=\scriptsize,left] at (0,0) {#2};
  \fill (1.8,1.6) circle (1.5pt);
  \ifnum#5=1
    \node[sw,draw=green!55!black,line width=1.2pt] (q) at (1.8,2.45) {#3};
    \draw[line width=0.8pt] (1.8,3.2) -- (q.north);
    \draw[line width=0.8pt] (q.south) -- (1.8,1.6);
    \coordinate (da) at (1.8,0.52); \coordinate (dk) at (1.8,1.08);
    \draw[line width=0.8pt] (1.8,0) -- (da); \draw[line width=0.8pt] (dk) -- (1.8,1.6);
  \else
    \node[sw,draw=green!55!black,line width=1.2pt] (q) at (1.8,0.75) {#3};
    \draw[line width=0.8pt] (1.8,1.6) -- (q.north);
    \draw[line width=0.8pt] (q.south) -- (1.8,0);
    \coordinate (da) at (1.8,2.12); \coordinate (dk) at (1.8,2.68);
    \draw[line width=0.8pt] (1.8,1.6) -- (da); \draw[line width=0.8pt] (dk) -- (1.8,3.2);
  \fi
  % diyot: anot (da) altta, katot (dk) üstte
  \filldraw[fill=white,draw=orange!85!black,line width=1.1pt] ($(da)+(-0.24,0)$) -- ($(da)+(0.24,0)$) -- (dk) -- cycle;
  \draw[orange!85!black,line width=1.3pt] ($(dk)+(-0.24,0)$) -- ($(dk)+(0.24,0)$);
  \node[font=\scriptsize,left=1.2mm,align=right] at ($(da)!0.5!(dk)+(-0.24,0)$) {#4};
  \coordinate (l1a) at (2.45,1.6);
  \coordinate (l1b) at (3.35,1.6);
  \coordinate (cnode) at (3.88,1.6);
  \coordinate (l2a) at (4.30,1.6);
  \coordinate (l2b) at (5.20,1.6);
  \draw[line width=0.8pt] (1.8,1.6) -- (l1a);
  \draw[coilpath,line width=0.9pt] (l1a) -- (l1b);
  \draw[line width=0.8pt] (l1b) -- (cnode) -- (l2a);
  \draw[coilpath,line width=0.9pt] (l2a) -- (l2b);
  \node[font=\scriptsize,above=2mm] at ($(l1a)!0.5!(l1b)$) {$L_1$};
  \node[font=\scriptsize,above=2mm] at ($(l2a)!0.5!(l2b)$) {$L_2$};
  \fill (cnode) circle (1.3pt);
  \draw[line width=0.8pt] (cnode) -- (3.88,1.0);
  \draw[line width=0.8pt] (3.55,1.0) -- (4.21,1.0);
  \draw[line width=0.8pt] (3.55,0.75) -- (4.21,0.75);
  \draw[line width=0.8pt] (3.88,0.75) -- (3.88,0.2);
  \draw[line width=0.8pt] (3.58,0.2) -- (4.18,0.2);
  \draw[line width=0.8pt] (3.68,0.05) -- (4.08,0.05);
  \node[font=\scriptsize,right] at (4.18,0.88) {$C_f$};
  \ifnum#5=1
    \draw[current] (l2b) -- (6.0,1.6);
    \draw[blue!70!black,line width=1.8pt] (0.5,3.2) -| (q.north);
    \draw[blue!70!black,line width=1.8pt] (q.south) -- (1.8,1.6) -- (l1a);
    \draw[freewheel] (0.5,0) -- (1.8,0) -- (da);
    \draw[freewheel] (dk) -- (1.8,1.6);
  \else
    \draw[current] (6.0,1.6) -- (l2b);
    \draw[blue!70!black,line width=1.8pt] (l1a) -- (1.8,1.6) -- (q.north);
    \draw[blue!70!black,line width=1.8pt] (q.south) -- (1.8,0) -- (0.5,0);
    \draw[freewheel] (1.8,1.6) -- (da);
    \draw[freewheel] (dk) -- (1.8,3.2) -- (0.5,3.2);
  \fi
  \node[font=\scriptsize,right] at (6.0,1.6) {$v_g,\ i_L$};
\end{tikzpicture}%
}


Literatür incelemesi, T-tipi dönüştürücünün yüksek verim ve güç yoğunluğu bakımından güçlü bir seçenek olduğunu göstermiştir. Aynı inceleme, hafif yükteki ters akımın, nötr kol kullanımının ve şebeke akımı harmoniklerinin aynı sistem sınırı içinde ele alınması gerektiğini de ortaya koymuştur. Tezde incelenen yapı, bölünmüş doğru akım (DA) bara, T-tipi güç katı, LCL çıkış filtresi, şebeke bağlantısı, ölçüm zinciri ve sayısal denetim biriminden oluşmaktadır. Kapsam iki düzeyde ayrılmaktadır. Laboratuvar prototipi üç fazlı, dört telli bir güç katıdır ve nötr hattı DA bara orta noktasına bağlıdır. Benzetim modelleri ile laboratuvar deneyleri ise prototipin tek faz kolunu kapsamaktadır. Şebeke dönüş akımı DA bara orta noktasına girmektedir. Bütün güç, akım ve verim değerleri faz başına verilmiştir. Deneyler güvenli laboratuvar koşulunda, varyak üzerinden yaklaşık \SI{120}{\volt} şebeke geriliminde ve yaklaşık \SI{2.3}{\ampere} faz akımında yürütülmüştür. Bölümün ilk amacı, önerilen yöntemlerin uygulanacağı fiziksel sistemi tanımlamak ve anahtarlama durumlarını ortak bir gösterim altında toplamaktır. Bölüm ayrıca sonraki analizlerin dayanacağı matematiksel modeli kurmakta ve benzetim modelinin tez tanımına göre nasıl doğrulandığını belgelemektedir.

Model, ideal anahtarlama durumlarını açıklayan devre düzeyini, LCL filtre ile denetim çevrimini açıklayan ortalama değer düzeyinden ayırmaktadır. Topolojiye ilişkin iddialar çok seviyeli dönüştürücü ve hafif yük literatürüyle sınanmıştır \citep{Mohan2003,Holmes2003,Schweizer2013,Wang2020,Chen2010}. Yarı iletken gerilim düşümleri, ters toparlanma, ölü zaman ve modülasyon doygunluğu anahtarlamalı benzetim katmanında ele alınmaktadır. Örnekleme gecikmesi ise kararlılık analizinde ayrıca ele alınmıştır.

\section{Sistemin Genel Mimarisi ve İnceleme Sınırı}

Şekil~\ref{fig:genel-sistem-mimarisi}, enerji, ölçüm ve denetim yollarını aynı çerçevede göstermektedir. Deşarj yönünde enerjinin kaynağı batarya veya ayarlanabilir DA kaynaktır. Enerji, bölünmüş DA bara üzerinden T-tipi güç katına aktarılmakta ve anahtarlanan faz gerilimleri LCL filtre aracılığıyla şebeke akımına dönüştürülmektedir. Akım ve gerilim ölçümleri sayısal denetim birimine taşınmakta, üretilen modülasyon komutları kapı sürücüleri üzerinden güç katına uygulanmaktadır. Sin\textsuperscript{3} kompanzasyonunun hedefi, nötr anahtar susturma yönteminin düşük güçte oluşturduğu üçüncü harmonik bileşendir.

\begin{figure}[H]
  \centering
  \resizebox{0.98\textwidth}{!}{%
  \begin{tikzpicture}[
    node distance=8mm and 8mm,
    block/.style={draw, rounded corners=1.5pt, align=center, minimum height=11mm, minimum width=24mm, fill=blue!5},
    ctrl/.style={draw, rounded corners=1.5pt, align=center, minimum height=10mm, minimum width=25mm, fill=orange!10},
    meas/.style={draw, rounded corners=1.5pt, align=center, minimum height=9mm, minimum width=24mm, fill=gray!10},
    flow/.style={-{Latex[length=2mm]}, thick},
    sig/.style={-{Latex[length=2mm]}, thick, dashed}
  ]
    \node[block] (bat) {Batarya /\\DA kaynak};
    \node[block, right=of bat] (dc) {Bölünmüş\\DA bara};
    \node[block, right=of dc] (inv) {T-tipi\\güç katı};
    \node[block, right=of inv] (lcl) {LCL\\filtre};
    \node[block, right=of lcl] (grid) {Şebeke /\\ortak bağlantı};
    \draw[flow] (bat) -- node[above, font=\scriptsize]{enerji} (dc);
    \draw[flow] (dc) -- (inv);
    \draw[flow] (inv) -- (lcl);
    \draw[flow] (lcl) -- (grid);

    \node[meas, below=15mm of lcl] (meas) {Gerilim ve akım\\ölçümleri};
    \node[ctrl, left=of meas] (pwm) {PWM ve kapı\\sinyalleri};
    \node[ctrl, left=of pwm] (control) {Akım denetimi ve\\sınırlandırıcılar};
    \node[ctrl, below=of control] (sin3) {{Sin\textsuperscript{3}}\\kompanzasyonu};
    \draw[sig] (grid.south) |- (meas.east);
    \draw[sig] (lcl.south) -- (meas.north);
    \draw[sig] (meas) -- (control);
    \draw[sig] (control) -- (pwm);
    \draw[sig] (sin3) -- (control);
    \draw[sig] (pwm.north) |- (inv.south);
  \end{tikzpicture}%
  }
  \caption{Genel sistem mimarisi: güç yolu, geri besleme yolu ve Sin\textsuperscript{3} kompanzasyonu.}
  \label{fig:genel-sistem-mimarisi}
\end{figure}

İnceleme sınırı üç katman üzerinden kurulmuştur. Fiziksel katman yarı iletkenleri, bölünmüş DA barayı, filtreyi ve şebeke bağlantısını kapsamaktadır. Denetim katmanı ölçüm, akım regülasyonu, modülasyon ve kapı sinyali üretimini içermektedir. Doğrulama katmanının kapsamı ise benzetim ile deneyde ortak kullanılan çalışma noktaları ve performans ölçütleridir. GT-tipi yeniden yapılandırma, aynı yarı iletkenlerin şarj ve deşarj kiplerinde farklı görevlerde kullanılmasına dayanmaktadır. Ayrıntılı işlev paylaşımı Bölüm~\ref{chap:gttipi} altında ele alınmıştır. Bölüm, deşarj yönündeki T-tipi evirici davranışını ve önerilen düşük güç yöntemlerinin dayandığı modeli sabitlemektedir.

Her iddia, kendisini en doğrudan sınayan model katmanıyla desteklenmiştir. Tablo~\ref{tab:model-katmanlari}, bölümde kullanılan model katmanlarını, her katmanın kapsadığı etkileri, sınadığı konuyu ve tamamlayıcı katmanını birlikte vermektedir. Katman ayrımı, ortalama değerli denetim modelinden elde edilen kararlılık sonucunu anahtarlama kaybı ve deneysel verim kanıtından ayrı tutmaktadır.

\begin{table}[H]
  \centering
  \caption{Model ve doğrulama katmanlarının amacı, sınadığı konu ve tamamlayıcı katmanı.}
  \label{tab:model-katmanlari}
  \small
  \begin{tabularx}{\textwidth}{p{0.21\textwidth} p{0.27\textwidth} X X}
    \toprule
    Katman & Temel içerik & Sınanan konu & Tamamlayıcı katman ve konu \\
    \midrule
    İdeal anahtarlama durumu & $P$--$O$--$N$ seviyeleri ve akım yönü & Geçerli kapı durumu, gerilim seviyesi ve iletim yolu & Anahtarlamalı benzetim: ölü zaman, kayıp ve THD \\
    Ortalama LCL/denetim modeli & $L_1$--$C_f$--$L_2$, gecikme ve $C_i(s)$ & Akım izleme, kutuplar ve kararlılık marjı & Anahtarlamalı benzetim: yarı iletken geçiş kaybı \\
    Anahtarlamalı benzetim & Gerçek PWM, kapı mantığı ve ideal olmayan elemanlar & Ters akım, harmonikler, anahtarlama sayısı ve geçici rejim & Deneysel sistem: donanımda ölçülen verim ve THD \\
    Deneysel sistem & Güç katı, sensör, sürücü, ADC ve gerçek zamanlı kod & Verim, THD, koruma ve uygulanabilirlik & S4 benzetimi: orta ve anma akımı bölgesi \\
    \bottomrule
  \end{tabularx}
\end{table}

\section{T-Tipi Faz Kolu ve Anahtarlama Durumları}

Faz kolunun devre yapısı ve $Q_1$--$Q_4$ adlandırması Şekil~\ref{fig:t-tipi-faz-kolu} üzerinde verilmiştir. $Q_1$--$Q_4$, Tablo~\ref{tab:sistem-parametreleri} içindeki $S_1$--$S_4$ anahtarlarıyla birebir aynıdır. Bölünmüş bara geriliminin toplamı $V_{\mathrm{dc}}$ kabul edildiğinde ideal faz-kutup gerilimi $+V_{\mathrm{dc}}/2$, $0$ ve $-V_{\mathrm{dc}}/2$ seviyelerinden birini almaktadır. Nötr kolun eklediği sıfır seviyesi, iki seviyeli yapıya göre daha küçük gerilim basamakları oluşturmaktadır. Karşılığında komütasyon düzeninin ve nötr nokta akımının ayrıca yönetilmesi gerekmektedir \citep{Schweizer2013,Wang2020}.

Sistemin tam açık güç şeması Şekil~\ref{fig:gelen-tam-topoloji} üzerinde, birinci faz hücresi ise okunabilir ölçekte Şekil~\ref{fig:gelen-faz1-hucresi} üzerinde verilmiştir. Sistem sınırı şemaya göre belirlenmiştir. Şemanın solunda dört iletkenli giriş, ortak-mod endüktansları ve şönt kapasitörlerden oluşan giriş/EMI ağı yer almaktadır. Ortada, ortak DA bağlantısına bağlanan üç tekrarlı T-tipi faz hücresi bulunmaktadır. Her hücrede dış kol anahtar modülü, çift yönlü nötr kol, PFC yükseltici (boost) diyotları, yerel kapasitörler, akım algılama elemanı ve röle vardır. Hücrelerin altındaki nötr hattı, üç fazın DA bara orta noktalarını birbirine ve giriş nötrüne bağlamaktadır. Şemanın en alt kısmında ise GT-tipi mimarinin şarj yönünde kullanılan batarya yönetim sistemi (BYS) transformatörleri görülmektedir. Deney düzeneği GT-tipi mimariye göre kurulmuştur. Her hücredeki röle ($RL1$) ile doğrultucu ve yükseltici diyotlar, şarj ve deşarj kipleri arasındaki geçişe hizmet etmektedir. Şemadaki $L1$, $L2$ ve $L3$ etiketleri LCL filtrenin ardışık $L_1$--$L_2$ elemanlarını değil, sırasıyla A, B ve C fazlarına ait üç eş \SI{600}{\micro\henry} bobini göstermektedir. Dolayısıyla $L2=\SI{600}{\micro\henry}$ etiketi filtrenin şebeke tarafı endüktansı değildir. Tezde bütün topoloji üç amaçla kullanılmaktadır: sistem ve ölçüm sınırını belirlemek, üç fazın aynı faz-başına modelin tekrarı olduğunu göstermek ve yarı iletken/parazit parametrelerini gerçek devre elemanlarıyla eşlemek. Anahtarlama denklemleri ve akım yolu çizimleri okunabilirlik için Şekil~\ref{fig:t-tipi-akim-yollari} altındaki tek faz eşdeğerine indirgenmiştir. Tek faz indirgemesi yalnız gösterimi sadeleştirmektedir. Giriş filtresi, diğer iki faz ve ölçüm/koruma elemanları fiziksel sistemde yer almaktadır.

\begin{landscape}
\begin{figure}[H]
  \centering
  \includegraphics[width=\linewidth,height=0.78\textheight,keepaspectratio]{umit_2026-11-09/tam_topoloji_notr_q1q3_q2q4.pdf}
  \caption{GT-tipi olarak kurulan deney düzeneğinin üç fazlı tam güç topolojisi: ortak giriş/EMI ağı, üç tekrarlı T-tipi faz hücresi ve DA bara orta noktalarını giriş nötrüne bağlayan nötr hattı. Şema ayrıca akım algılama ile kip değişim rölesi yollarını ve şarj yönündeki BYS transformatörlerini içermektedir. Her hücrede nötr kolun faz tarafındaki anahtarı dış kolun üst anahtarıyla, orta nokta tarafındaki anahtarı alt anahtarıyla tümleyen çift oluşturmaktadır (birinci fazda $Q_1$--$Q_3$ ve $Q_2$--$Q_4$). Şema, yöntem bölümünde sistem sınırı ve faz-başına indirgeme kaynağı olarak kullanılmıştır.}
  \label{fig:gelen-tam-topoloji}
\end{figure}
\end{landscape}

\begin{landscape}
\begin{figure}[H]
  \centering
  \includegraphics[width=\linewidth,height=0.78\textheight,keepaspectratio]{umit_2026-11-09/tam_topoloji_faz1_hucresi.pdf}
  \caption{Şekil~\ref{fig:gelen-tam-topoloji} içindeki birinci faz hücresinin büyütülmüş görünümü. Hücrede $D_1$--$D_2$ PFC yükseltici (boost) diyotları, $C_{10}$--$C_{11}$ bölünmüş DA bara kapasitörleri, dış kol modülü $Q_1$--$Q_2$ (SKM100GB12T4) ve $C_{12}$ snubber kapasitörü yer almaktadır. Hücrede ayrıca orta nokta tarafında $Q_4$ ve faz tarafında $Q_3$ nötr IGBT'leri ile $RL1$--$RL2$ kip değişim röleleri bulunmaktadır. Alttaki yatay hat, faz hücrelerinin DA bara orta noktalarını birbirine bağlayan nötr hattıdır.}
  \label{fig:gelen-faz1-hucresi}
\end{figure}
\end{landscape}

Faz indisi $x\in\{a,b,c\}$, DA bara orta noktası $N$, faz kutup gerilimi $v_{xN}$ ve gerilim durumunu seçen boyutsuz anahtarlama değişkeni $s_x\in\{-1,0,1\}$ ile gösterilsin. Üç gerilim durumu
\begin{equation}
  v_{xN}=s_x\frac{V_{\mathrm{dc}}}{2}
  \label{eq:t-tipi-anahtarlama-fonksiyonu}
\end{equation}
biçiminde ortak bir gösterime indirgenebilir. Eşitlik~\eqref{eq:t-tipi-anahtarlama-fonksiyonu}, $s_x=1$, $0$ ve $-1$ için sırasıyla pozitif, sıfır ve negatif kutup gerilimini vermektedir. Doğrusal modülasyon bölgesinde normalize faz komutu $m_x=2v_x^{\ast}/V_{\mathrm{dc}}$ olmak üzere pozitif, negatif ve sıfır durumlarının boyutsuz ideal görev oranları $d_P$, $d_N$ ve $d_O$
\begin{equation}
  d_P=\max(m_x,0),\qquad
  d_N=\max(-m_x,0),\qquad
  d_O=1-|m_x|,\qquad |m_x|\leq 1
  \label{eq:t-tipi-gorev-oranlari}
\end{equation}
olarak yazılabilir \citep{Holmes2003,Wang2020}. Eşitlik~\eqref{eq:t-tipi-gorev-oranlari}, modülasyon komutunun genliği azaldığında sıfır durumunun payının neden arttığını doğrudan göstermektedir. Gerçek taşıyıcı tabanlı PWM'de örnekleme, minimum darbe, ölü zaman ve nötr nokta dengeleme müdahaleleri ideal görev oranlarını değiştirebilir. Benzetim modelinde ölü zaman $t_d$, her kapı komutunun yükselen kenarına uygulanan gecikmeyle temsil edilmektedir (Tablo~\ref{tab:sistem-parametreleri}). Görev oranları, ideal değerleriyle doğrudan PWM karşılaştırıcısına aktarılmaktadır. Ölü zamanın düşük akımdaki etkisi kalibrasyon benzetimlerinde ölçülmüştür (Alt Bölüm~\ref{sec:kalibrasyon}).

Üç gerilim seviyesinin devre karşılığı Şekil~\ref{fig:t-tipi-uc-seviye} üzerinde tek karşılaştırmalı görünüm altında verilmiştir. Aynı devre tabanının korunması, hangi anahtarlama durumunun yalnızca bağlantı noktasını değiştirdiğini görmeyi kolaylaştırmaktadır.

\begin{figure}[H]
  \centering
  \begin{subfigure}[t]{0.32\textwidth}
    \centering
    \resizebox{\linewidth}{!}{\ttypevector{green!55!black}{gray!65}{gray!65}{gray!65}{%
      \draw[activepath] (4,5) -- (q1.north);
      \draw[activepath] (q1.south) -- (out) -- (5.0,2);}}
    \caption{$P$: $+V_{\mathrm{dc}}/2$}
  \end{subfigure}\hfill
  \begin{subfigure}[t]{0.32\textwidth}
    \centering
    \resizebox{\linewidth}{!}{\ttypevector{gray!65}{gray!65}{green!55!black}{green!55!black}{%
      \draw[activepath] (0,2) -- (q3.west);
      \draw[activepath] (q3.east) -- (q4.west);
      \draw[activepath] (q4.east) -- (out) -- (5.0,2);}}
    \caption{$O$: $0$}
  \end{subfigure}\hfill
  \begin{subfigure}[t]{0.32\textwidth}
    \centering
    \resizebox{\linewidth}{!}{\ttypevector{gray!65}{green!55!black}{gray!65}{gray!65}{%
      \draw[activepath] (4,-1) -- (q2.south);
      \draw[activepath] (q2.north) -- (out) -- (5.0,2);}}
    \caption{$N$: $-V_{\mathrm{dc}}/2$}
  \end{subfigure}
  \caption{T-tipi faz kolunun üç gerilim durumu. Yeşil anahtarlar ve mavi bağlantılar ilgili durumun etkin elemanlarını göstermektedir.}
  \label{fig:t-tipi-uc-seviye}
\end{figure}

Tablo~\ref{tab:t-tipi-durumlar}, kapı durumlarını literatürdeki üç seviyeli gösterimle birlikte özetlemektedir \citep{Holmes2003,Wang2020}. ``1'' iletim komutu, ``0'' kesim komutudur. Tablonun kapsamı ideal kararlı durumlardır. Ölü zaman aralığındaki diyot iletimi ve akım yönüne bağlı geçici komütasyon yolları, anahtarlamalı benzetimde yarı iletken modelleri üzerinden kendiliğinden oluşmaktadır.

\begin{table}[H]
  \centering
  \caption{T-tipi faz kolunun ideal anahtarlama durumları ve kapı komutları.}
  \label{tab:t-tipi-durumlar}
  \small
  \begin{tabularx}{\textwidth}{c c c c c c X}
    \toprule
    Durum & $Q_1$ & $Q_2$ & $Q_3$ & $Q_4$ & $v_{xN}$ & İşlevsel açıklama \\
    \midrule
    $P$ & 1 & 0 & 0 & 1 & $+V_{\mathrm{dc}}/2$ & Pozitif yarım periyotta enerji üst DA bara kolundan faz çıkışına aktarılmaktadır. \\
    $O$ & 0 & 0 & 1 & 1 & $0$ & Faz çıkışı nötr kola bağlanmaktadır. Filtre akımının yönüne göre serbest dolaşım yolu oluşmaktadır. \\
    $N$ & 0 & 1 & 1 & 0 & $-V_{\mathrm{dc}}/2$ & Negatif yarım periyotta enerji alt DA bara kolu üzerinden aktarılmaktadır. \\
    \bottomrule
  \end{tabularx}
\end{table}

$P$ durumunda $Q_4$ ve $N$ durumunda $Q_3$ için verilen iletim komutu, yarım periyot etkinleştirme kuralını yansıtmaktadır. İlgili nötr anahtar, aktif enerji yolunda akım taşımamaktadır. Yalnızca gerilim seviyesini üreten asgari durum gösterilecekse $P$ için $Q_1$, $O$ için $Q_3$--$Q_4$ ve $N$ için $Q_2$ iletimde kabul edilebilir. Benzetim modelinin PWM alt sistemi yarım periyot etkinleştirme gösterimini kullanmaktadır. Klasik anahtarlamada $Q_3$, pozitif yarım periyotta $Q_1$ komutunun tümleyenidir ve $Q_4$ sürekli iletimdedir. Negatif yarım periyotta roller simetrik olarak $Q_3$ ile $Q_4$ arasında yer değiştirmektedir. Modelden okunan kapı mantığı Tablo~\ref{tab:t-tipi-durumlar} ile birebir örtüşmektedir. Susturmalı anahtarlamanın kapı mantığı Alt Bölüm~\ref{sec:model-kapi-mantigi} içinde verilmiştir.

Modülasyon, pozitif yarım periyotta $P$ ile $O$, negatif yarım periyotta ise $N$ ile $O$ durumları arasında yapılmaktadır. Endüktif çıkış filtresi nedeniyle faz akımı anahtarlama anında sürekliliğini korumaktadır. Gerilim komutu değişse bile akımın mevcut yönü, aktif anahtarın kanalı veya antiparalel diyot üzerinden uygun serbest dolaşım yolunu belirlemektedir. Dolayısıyla gerçek iletimdeki yarı iletken, kapı komutlarının yanında akım yönüne de bağlıdır.

Şekil~\ref{fig:t-tipi-akim-yollari}(a), $C_A$--$C_B$ bölünmüş DA barayı, antiparalel diyotlu $Q_1$--$Q_4$ IGBT'lerini ve faz başına $L_1$--$C_f$--$L_2$ filtresini tek topolojide birleştirmektedir. Şekil~\ref{fig:t-tipi-akim-yollari}(b)--(e), aynı güç katının dört temel iletim yolunu karşılaştırmaktadır. Akım yolunun okunabilirliğini korumak için filtre elemanları alt çizimlerde tekrarlanmamıştır. Pozitif alternansın aktif durumunda akış $Q_1$ üzerinden, sıfır durumunda nötr kol üzerinden gerçekleşmektedir. Negatif alternansta aynı işleyiş alt kol için simetriktir. Nötr seviyesinde geçirilen sürenin modülasyon indeksi azaldıkça artması, düşük güçte nötr anahtarların iletim payını ve ters akım riskini büyütmektedir.

\begin{figure}[H]
  \centering
  \begin{subfigure}[t]{0.96\textwidth}
    \centering
    \resizebox{0.98\linewidth}{!}{\ttypefulltopology}
    \caption{Faz başına tam güç topolojisi ve LCL filtre}
  \end{subfigure}

  \vspace{4mm}
  \begin{subfigure}[t]{0.48\textwidth}
    \centering
    \ttypevector{green!55!black}{gray!65}{gray!65}{gray!65}{%
      \draw[current] (0.7,5) -- (q1.north);
      \draw[current] (q1.south) -- (out) -- (5.0,2);}
    \caption{Pozitif alternans, aktif $P$ durumu}
  \end{subfigure}\hfill
  \begin{subfigure}[t]{0.48\textwidth}
    \centering
    \ttypevector{gray!65}{gray!65}{green!55!black}{green!55!black}{%
      \draw[current] (0.5,2) -- (q3.west);
      \draw[current] (q3.east) -- (q4.west);
      \draw[current] (q4.east) -- (5.0,2);}
    \caption{Pozitif alternans, sıfır $O$ durumu}
  \end{subfigure}

  \vspace{4mm}
  \begin{subfigure}[t]{0.48\textwidth}
    \centering
    \ttypevector{gray!65}{green!55!black}{gray!65}{gray!65}{%
      \draw[current] (5.0,2) -- (out) -- (q2.north);
      \draw[current] (q2.south) -- (0.7,-1);}
    \caption{Negatif alternans, aktif $N$ durumu}
  \end{subfigure}\hfill
  \begin{subfigure}[t]{0.48\textwidth}
    \centering
    \ttypevector{gray!65}{gray!65}{green!55!black}{green!55!black}{%
      \draw[current] (5.0,2) -- (q4.east);
      \draw[current] (q4.west) -- (q3.east);
      \draw[current] (q3.west) -- (0.5,2);}
    \caption{Negatif alternans, sıfır $O$ durumu}
  \end{subfigure}
  \caption{T-tipi faz kolunun devre yapısı ile dört temel akım yolunun ortak devre dili ve renk lejandıyla vektörel gösterimi. Tam topoloji $C_A$, $C_B$, antiparalel diyotlu IGBT'ler, $L_1$, $C_f$ ve $L_2$ elemanlarını içermektedir. Kırmızı oklar, tanımlanan pozitif akım yönüne göre iletim yollarını göstermektedir.}
  \label{fig:t-tipi-akim-yollari}
\end{figure}

\begin{samepage}
  Şekil~\ref{fig:t-tipi-akim-yollari}(a), faz kolunun işlevsel güç topolojisidir. Sensör, koruma ve giriş filtresi elemanlarının yeri Şekil~\ref{fig:gelen-tam-topoloji} üzerinde, yarı iletken parça kodları ve pasif eleman değerleri ise Tablo~\ref{tab:sistem-parametreleri} içinde verilmiştir.
\end{samepage}

\section{Düşük Güç Eşdeğer Modeli}
\label{sec:dusuk-guc-esdeger}

Faz gerilimi referansının genliği azaldığında aktif $P$ veya $N$ durumunun görev oranı küçülmekte, sıfır durumunun süresi uzamaktadır. Pozitif ve negatif yarım periyotlar ayrı ayrı ele alındığında faz kolu, anahtarlanan gerilim kaynağı ile filtre endüktansı bakımından senkron alçaltıcı dönüştürücüye eşdeğer bir davranış göstermektedir \citep{Mohan2003,Erickson2001}. İki yarım periyodun eşdeğerleri Şekil~\ref{fig:senkron-buck-esdeger} altında bir araya getirilmiştir.

\begin{figure}[H]
  \centering
  \begin{subfigure}[t]{0.46\textwidth}
    \centering
    \resizebox{\linewidth}{!}{\syncbuckvector{$+V_{\mathrm{dc}}/2$}{$N$}{$Q_1$}{$Q_3$}{1}}
    \caption{Pozitif yarım periyot eşdeğeri}
  \end{subfigure}\hfill
  \begin{subfigure}[t]{0.46\textwidth}
    \centering
    \resizebox{\linewidth}{!}{\syncbuckvector{$N$}{$-V_{\mathrm{dc}}/2$}{$Q_4$}{$Q_2$}{0}}
    \caption{Negatif yarım periyot eşdeğeri}
  \end{subfigure}
  \caption{T-tipi faz kolunun yarım periyotlara göre senkron-buck eşdeğerleri. Yeşil anahtar ve mavi yol aktif enerji aktarımını göstermektedir.}
  \label{fig:senkron-buck-esdeger}
\end{figure}

Anahtarlama periyodu $T_s$, boyutsuz görev oranı $d\in[0,1]$, faz koluna uygulanan etkin DA gerilimi $V_h$, anlık şebeke gerilimi $v_g$ ve eşdeğer filtre endüktansı $L$ olarak tanımlansın. Anahtarın iletim ve kesim aralıklarındaki endüktör gerilimleri de sırasıyla $v_{L,\mathrm{on}}$ ve $v_{L,\mathrm{off}}$ olsun. Pozitif yarım periyotta ideal gerilimler
\begin{equation}
  \begin{aligned}
    v_{L,\mathrm{on}} &= V_h-v_g, \\
    v_{L,\mathrm{off}} &= -v_g
  \end{aligned}
  \label{eq:buck-enduktor-gerilimleri}
\end{equation}
biçimindedir. Eşitlik~\eqref{eq:buck-enduktor-gerilimleri} altında anahtarlama periyodu boyunca gerilimlerin yaklaşık sabit kaldığı ve yerel kararlı durumda volt--saniye dengesinin sağlandığı kabul edilirse tepe--tepe akım dalgalanması $\Delta i_L$
\begin{equation}
  \Delta i_L \simeq \frac{(V_h-v_g)dT_s}{L}
  =\frac{v_g(1-d)T_s}{L}
  \label{eq:buck-akim-dalgalanmasi}
\end{equation}
ile bulunmaktadır \citep{Erickson2001,Chen2010}. Eşitlik~\eqref{eq:buck-akim-dalgalanmasi} içindeki iki ifade yalnızca aynı yerel volt--saniye dengesi altında eşdeğerdir. Ortalama akım $\bar{i}_L$ değerinin dalgalanmanın yarısından küçük olduğu bölgede akım sıfırı geçebilir:
\begin{equation}
  \bar{i}_L < \frac{\Delta i_L}{2}.
  \label{eq:ters-akim-esigi}
\end{equation}
Diyotlu alçaltıcı dönüştürücüde akım sıfıra ulaştığında serbest dolaşım diyodu doğal olarak kesime gitmektedir. Senkron yapıda diyoda paralel anahtar iletimde tutulduğu için doğal kesim yerine $i_L<0$ aralığı oluşmakta ve filtredeki enerji DA tarafa geri taşınmaktadır. Diyot benzetimi, akım sıfıra yaklaştığında senkron yolu keserek negatif akımı engellemektedir. Sıfır akım algılamadaki gecikme ve anahtarlama geçişleri ise gerçek kayıp sınırını belirleyen etkenlerdir \citep{Chen2010,Erickson2001}. İki kipin farkı Şekil~\ref{fig:de-ccm-akim} üzerinde gösterilmiştir.

Ortalama akım ile akım dalgalanması arasındaki ilişki üç çalışma bölgesi tanımlamaktadır \citep{Mohan2003,Erickson2001,Holmes2003}:

\begin{enumerate}
  \item \textbf{Sürekli iletim bölgesi:} $|\bar i_L|>\Delta i_L/2$ koşulunda akım bir anahtarlama periyodu içinde sıfıra ulaşmamaktadır. Senkron nötr anahtarın iletim gerilim düşümünü azaltması verim açısından yararlıdır.
  \item \textbf{Sınır iletim bölgesi:} $|\bar i_L|\simeq\Delta i_L/2$ noktasında akım periyot sonunda sıfıra yaklaşmaktadır. Parametre belirsizlikleri ve gecikmeler nedeniyle yöntem geçişi için histerezis gereklidir.
  \item \textbf{Zorlanmış sürekli iletim bölgesi:} $|\bar i_L|<\Delta i_L/2$ olmasına rağmen senkron anahtarlama sürdürülürse akım sıfırın ötesine geçmektedir. Ters akım kaybı belirginleşirken akımın sürekliliği harmonik performansı destekleyebilir.
\end{enumerate}

Bölgeler arasındaki sınır akımı
\begin{equation}
  I_{\mathrm{B}}(t)=\frac{\Delta i_L(t)}{2}
  \label{eq:sinir-akimi}
\end{equation}
ile tanımlanmaktadır. $I_{\mathrm{B}}$, DA bara gerilimine, anlık şebeke gerilimine, görev oranına, anahtarlama frekansına ve eşdeğer endüktansa bağlıdır. Sabit bir yüzde güç eşiği yerine fiziksel değişkenlerden türetilen sınır akımı, şebeke gerilimi veya DA bara seviyesi değiştiğinde daha tutarlı bir kip geçişi sağlamaktadır.

\section{LCL Filtre ve Şebeke Modeli}

Şebeke bağlantılı eviricide LCL filtre, anahtarlama frekansı çevresindeki gerilim bileşenlerini daha küçük toplam endüktansla bastırabilmektedir. Rezonans tepesi ve şebeke empedansına duyarlılık, filtreyi yalnızca pasif bir zayıflatıcı olmaktan çıkarıp kapalı çevrim tasarımının parçası hâline getirmektedir \citep{Zhang2014PR,Lu2016}. Faz başına ortalama modelde evirici tarafı endüktansı $L_1$, filtre kapasitörü $C_f$, şebeke tarafı endüktansı $L_2$, seri kayıplar $R_1$ ve $R_2$ ile gösterilmiştir. Şebeke eşdeğer endüktansı gerektiğinde $L_2$ içine eklenmektedir.

Şekil~\ref{fig:faz-basina-lcl-modeli}, durum denklemlerinde kullanılan işaretleri fiziksel ölçüm noktalarıyla eşleştirmektedir. $i_1$ evirici tarafı endüktör akımı, $i_2$ ortak bağlantı noktasına aktarılan şebeke akımı, $v_c$ ise filtre kapasitörü gerilimidir. Ölçülen akımın $i_1$ veya $i_2$ olması, denetlenen sistemin modelini ve aktif sönümleme yolunu değiştirmektedir. Benzetim modelinde güç, RMS ve THD ölçümleri EMC süzgecinden sonra, $L_2$'den önce, evirici çıkış uçlarında (C noktası) alınmaktadır. Ölçülen akım, filtre kapasitörü ile X kapasitörünün akımını içeren şebeke akımıdır. S4 akım denetimi ise $i_1=i_{L1}$ akımını geri beslemektedir (Tablo~\ref{tab:sistem-parametreleri}).

\begin{figure}[H]
  \centering
  \resizebox{0.96\textwidth}{!}{%
  \begin{tikzpicture}[
    block/.style={draw,rounded corners=1pt,minimum width=13mm,minimum height=8mm,fill=blue!4},
    resistor/.style={draw,minimum width=8mm,minimum height=5mm,fill=white},
    coilpath/.style={decorate,decoration={coil,aspect=0.48,segment length=2.9mm,amplitude=1.9mm}},
    sensor/.style={circle,draw,fill=orange!10,minimum size=7mm,font=\scriptsize},
    flow/.style={-{Latex[length=2mm]},line width=1.0pt},
    current/.style={-{Latex[length=2mm]},red!80!black,line width=1.0pt},
    wire/.style={line width=0.9pt},
    signal/.style={-{Latex[length=1.8mm]},dashed,gray!80}
  ]
    \node[block,align=center] (inv) at (0,0) {T-tipi\\faz kolu};
    \node[resistor] (r1) at (1.65,0) {$R_1$};
    \coordinate (l1a) at (2.35,0);
    \coordinate (l1b) at (3.45,0);
    \coordinate (vc) at (4.10,0);
    \node[resistor] (r2) at (5.30,0) {$R_2$};
    \coordinate (l2a) at (5.95,0);
    \coordinate (l2b) at (7.05,0);
    \node[block,align=center] (grid) at (8.55,0) {Şebeke\\$v_g$};
    \draw[wire] (inv.east) -- node[above,font=\scriptsize] {$v_{\mathrm{inv}}$} (r1.west);
    \draw[wire] (r1.east) -- (l1a);
    \draw[coilpath,line width=0.9pt] (l1a) -- (l1b);
    \draw[wire] (l1b) -- (vc) -- (r2.west);
    \draw[wire] (r2.east) -- (l2a);
    \draw[coilpath,line width=0.9pt] (l2a) -- (l2b);
    \draw[wire] (l2b) -- (grid.west);
    \node[above=2mm,font=\scriptsize] at ($(l1a)!0.5!(l1b)$) {$L_1$};
    \node[above=2mm,font=\scriptsize] at ($(l2a)!0.5!(l2b)$) {$L_2$};
    \draw[current] (2.30,0.55) -- node[above,font=\scriptsize] {$i_1$} (3.45,0.55);
    \draw[current] (5.90,0.55) -- node[above,font=\scriptsize] {$i_2$} (7.05,0.55);
    \fill (vc) circle (1.6pt);
    \draw[line width=0.8pt] (vc) -- (4.10,-0.8);
    \draw[line width=0.8pt] (3.73,-0.8) -- (4.47,-0.8);
    \draw[line width=0.8pt] (3.73,-1.08) -- (4.47,-1.08);
    \draw[line width=0.8pt] (4.10,-1.08) -- (4.10,-1.50);
    \draw[line width=0.8pt] (3.73,-1.50) -- (4.47,-1.50);
    \draw[line width=0.8pt] (3.83,-1.65) -- (4.37,-1.65);
    \draw[line width=0.8pt] (3.95,-1.80) -- (4.25,-1.80);
    \node[right,font=\scriptsize] at (4.40,-0.94) {$C_f$};
    \node[above,font=\scriptsize] at (vc) {$v_c$};
    \node[sensor] (si1) at (2.90,-2.25) {$i_1$};
    \node[sensor] (svc) at (4.10,-2.25) {$v_c$};
    \node[sensor] (si2) at (6.50,-2.25) {$i_2$};
    \node[block,fill=gray!8,align=center] (adc) at (4.85,-3.55) {Ölçüm ve\\ADC zinciri};
    \draw[signal] ($(l1a)!0.5!(l1b)$) -- (si1.north);
    \draw[signal] (vc) -- (svc.north);
    \draw[signal] ($(l2a)!0.5!(l2b)$) -- (si2.north);
    \draw[signal] (si1) -- (adc);
    \draw[signal] (svc) -- (adc);
    \draw[signal] (si2) -- (adc);
  \end{tikzpicture}%
  }
  \caption{Faz başına LCL filtre, durum değişkenleri ve denetimde kullanılabilecek ölçüm noktaları.}
  \label{fig:faz-basina-lcl-modeli}
\end{figure}

Evirici kutup gerilimi $v_{\mathrm{inv}}$, kapasitör gerilimi $v_c$, şebeke gerilimi $v_g$, evirici tarafı akımı $i_1$ ve şebeke akımı $i_2$ için durum denklemleri
\begin{equation}
  \begin{aligned}
    L_1\frac{\mathrm{d}i_1}{\mathrm{d}t} &= v_{\mathrm{inv}}-v_c-R_1i_1, \\
    C_f\frac{\mathrm{d}v_c}{\mathrm{d}t} &= i_1-i_2, \\
    L_2\frac{\mathrm{d}i_2}{\mathrm{d}t} &= v_c-v_g-R_2i_2
  \end{aligned}
  \label{eq:lcl-durum-denklemleri}
\end{equation}
olarak elde edilmektedir. Eşitlik~\eqref{eq:lcl-durum-denklemleri} içindeki $L_1$, $L_2$ ve $C_f$ enerji depolayan filtre elemanlarıdır. $R_1$ ile $R_2$ ise seri eşdeğer kayıpları temsil etmektedir. Laplace değişkeni $s$ olmak üzere $i_1(t)$ ve $v_{\mathrm{inv}}(t)$ işaretlerinin Laplace dönüşümleri sırasıyla $I_1(s)$ ve $V_{\mathrm{inv}}(s)$ olarak tanımlansın. Kayıpların ihmal edildiği modelde $v_g=0$ alınarak evirici kutup geriliminden evirici tarafı akımına aktarım
\begin{equation}
  G_{vi}(s)=\frac{I_1(s)}{V_{\mathrm{inv}}(s)}
  =\frac{L_2C_fs^2+1}{L_1L_2C_fs^3+(L_1+L_2)s}
  \label{eq:lcl-aktarim}
\end{equation}
şeklindedir. Eşitlik~\eqref{eq:lcl-aktarim} kayıpsız ve şebeke gerilimi bozucusu sıfırlanmış ideal denetlenen sistemi tanımlamaktadır. Rezonans açısal frekansı $\omega_r$ ve hertz cinsinden rezonans frekansı $f_r$
\begin{equation}
  \omega_r=\sqrt{\frac{L_1+L_2}{L_1L_2C_f}},
  \qquad f_r=\frac{\omega_r}{2\pi}
  \label{eq:lcl-rezonans}
\end{equation}
ile hesaplanmaktadır. Eşitlik~\eqref{eq:lcl-rezonans}, seri kayıpları ve ek sönümleme kolunu içermeyen ideal LCL sonucudur. Gerçek modelde seri dirençler, kapasitör sönümleme kolu, sayısal gecikme ve şebeke empedansı aktarım fonksiyonuna eklenmelidir. Kararlılık incelemesinde kullanılan $L(s)$, $S(s)$ ve $T(s)$ fonksiyonları, EMC süzgecini, ölçüm süzgecini ve denetim gecikmesini içeren fiziksel model üzerinden Alt Bölüm~\ref{sec:kararlilik-bode} içinde üretilmiştir.

\section{Denetim, Modülasyon ve Sin\texorpdfstring{\textsuperscript{3}}{3} Kompanzasyonu}

Akım çevrimi, referans akım $i^{\ast}$ ile ölçülen akım arasındaki hatadan modülasyon komutu üretmektedir. Laboratuvar deneyleri (S1--S3) açık çevrim modülasyon referansıyla yürütülmüştür. S4 benzetiminde akım denetleyicisi $C_i(s)$, $L_1$ akımını geri besleyen oransal (P), oransal--integral (PI) ya da oransal--integral--rezonans (PIR) yapıdadır \citep{Teodorescu2006PR,Zhang2014PR}. PIR yapısının rezonans terimi şebeke temel açısal frekansı $\omega_0$ çevresinde kazanç sağlamakta, rezonans bandını ise $\omega_b$ belirlemektedir. Denetleyici tanımı ve tasarımı Alt Bölüm~\ref{sec:s4-kapali-cevrim} içinde, katsayılar Tablo~\ref{tab:sistem-parametreleri} içinde verilmiştir. PI sıfırı, akım çevriminin yaklaşık \SI{586}{\hertz} olan kazanç geçiş frekansının onda biri düzeyinde, \SI{60}{\hertz}'e yerleştirilmiştir. Seçim, şebeke frekansından bağımsız olarak kazanç geçiş frekansına dayanmaktadır.

Referans ve ölçülen akımların Laplace dönüşümleri $I^{\ast}(s)$ ile $I(s)$ olsun. Normalize toplam modülasyon komutu ve Sin\textsuperscript{3} ileri besleme yolu da sırasıyla $U_{\mathrm{mod}}^{\ast}(s)$ ve $U_{\sin^3}(s)$ olarak tanımlansın. Modülasyon komutu kavramsal olarak
\begin{equation}
  U_{\mathrm{mod}}^{\ast}(s)=C_i(s)\left[I^{\ast}(s)-I(s)\right]+U_{\sin^3}(s)
  \label{eq:modulasyon-toplami}
\end{equation}
şeklinde ayrılmıştır. Eşitlik~\eqref{eq:modulasyon-toplami} zaman alanı işaretleriyle aktarım fonksiyonunu aynı ifadede karıştırmamak için tamamen Laplace alanında yazılmıştır. PLL faz açısı $\theta$ için trigonometrik özdeşlik
\begin{equation}
  \sin^3(\theta)=\frac{3\sin(\theta)-\sin(3\theta)}{4}
  \label{eq:sin3-ozdeslik}
\end{equation}
kompanzasyonun temel ve üçüncü harmonik bileşenlerini açıkça göstermektedir. Eşitlik~\eqref{eq:sin3-ozdeslik} içindeki $\theta$ boyutsuz elektriksel açıdır. Benzetim modelinde terim $K_{3x}\sin^3\theta$ biçimindedir. $K_{1x}$ ve $K_{3x}$, Bölüm~\ref{chap:sin3} içindeki $k_1$ ve $k_3$ katsayılarının $x$ numaralı senaryodaki değerleridir. Birim genlikli referans sinüsünün küpü, pozitif işaretle ve ayrı bir normalizasyon olmadan görev oranı toplamına eklenmektedir. Eşitlik~\eqref{eq:sin3-ozdeslik} uygulandığında modülasyon referansı $(K_1+3K_3/4)\sin\theta-(K_3/4)\sin3\theta$ olmaktadır. Üçüncü harmonik bileşen temel bileşene ters işaretlidir ve tepe değer $K_{1x}+K_{3x}$'tür.

Ölçümden kapı işaretine uzanan işlevsel zincir Şekil~\ref{fig:denetim-modulasyon-blok} üzerinde izlenebilir bir blok yapısına dönüştürülmüştür. Nötr anahtar susturma kararı PWM üretiminden sonra kapı mantığında uygulanmaktadır.

\begin{figure}[H]
  \centering
  \resizebox{0.99\textwidth}{!}{%
  \begin{tikzpicture}[
    node distance=5mm and 5mm,
    block/.style={draw,rounded corners=1.3pt,align=center,minimum height=9mm,minimum width=14mm,fill=blue!5},
    special/.style={draw,rounded corners=1.3pt,align=center,minimum height=9mm,minimum width=18mm,fill=orange!12},
    sum/.style={circle,draw,minimum size=7mm,inner sep=0pt},
    flow/.style={-{Latex[length=1.8mm]},line width=0.9pt},
    feedback/.style={-{Latex[length=1.8mm]},line width=0.9pt,dashed}
  ]
    \node[block] (iref) {$i^{\ast}$\\üretimi};
    \node[sum,right=of iref] (sumi) {$\Sigma$};
    \node[block,right=of sumi] (ci) {$C_i(s)$\\akım denetimi};
    \node[sum,right=of ci] (sumv) {$\Sigma$};
    \node[block,right=of sumv,minimum width=24mm] (pwm) {Sınırlandırma, ölçekleme\\ve PWM üretimi};
    \node[special,right=of pwm,minimum width=25mm] (gate) {$Q_1$--$Q_4$ kapı mantığı\\ve nötr susturma};
    \node[block,right=of gate,minimum width=20mm] (plant) {T-tipi güç katı\\ve LCL filtre};
    \draw[flow] (iref) -- (sumi);
    \draw[flow] (sumi) -- node[above,font=\scriptsize] {$e_i$} (ci);
    \draw[flow] (ci) -- (sumv);
    \draw[flow] (sumv) -- node[above,font=\scriptsize] {$v_{\mathrm{mod}}^{\ast}$} (pwm);
    \draw[flow] (pwm) -- (gate);
    \draw[flow] (gate) -- node[above,font=\scriptsize] {$q_1$--$q_4$} (plant);

    \node[special,below=13mm of sumv,xshift=11mm,minimum width=20mm] (sin3) {{Sin\textsuperscript{3}}\\$K_{\sin^3}\sin^3\theta$};
    \node[block,below=11mm of iref] (pll) {Faz bilgisi\\$\theta$};
    \node[block,below=11mm of plant] (meas) {Ölçüm, ADC\\ve ölçekleme};
    \draw[flow] (pll) -- (iref);
    \coordinate (phasebus) at ($(pll.south)+(0,-0.65)$);
    \draw[flow] (pll.south) -- (phasebus) -| (sin3.south);
    \draw[flow] (sin3) -- (sumv.south);
    \draw[feedback] (plant.south) -- (meas.north);
    \coordinate (feedbackbus) at ($(sin3.south)+(0,-1.15)$);
    \draw[feedback] (meas.south) |- (feedbackbus) -| node[pos=0.72,below,font=\scriptsize] {$i_{\mathrm{ölç}}$} (sumi.south);
    \node[font=\scriptsize] at ($(sumi)+(-0.24,0.34)$) {$+$};
    \node[font=\scriptsize] at ($(sumi)+(-0.24,-0.34)$) {$-$};
    \node[font=\scriptsize] at ($(sumv)+(-0.24,0.34)$) {$+$};
    \pgfresetboundingbox
    \coordinate (bboxse) at (meas.east |- feedbackbus);
    \path[use as bounding box] ($(iref.north west)+(-0.3,0.35)$) rectangle ($(bboxse)+(0.3,-0.3)$);
  \end{tikzpicture}%
  }
  \caption{İşlevsel akım denetimi, Sin\textsuperscript{3} ileri besleme, PWM ve nötr anahtar susturma blokları.}
  \label{fig:denetim-modulasyon-blok}
\end{figure}

Sin\textsuperscript{3} ileri besleme terimi hatadan üretilmemekte ve geri besleme yolunda bulunmamaktadır. Dolayısıyla terim, denetlenen sistemin aktarım fonksiyonu $G(s)$ olmak üzere doğrusal modelde $1+C_i(s)G(s)=0$ karakteristik denklemini değiştirmemektedir (Alt Bölüm~\ref{sec:kararlilik-bode}). Sayısal uygulamada örnekleme, sıfırıncı derece tutucu, PWM, hesaplama ve ölçüm filtreleri toplam bir gecikme $G_d(s)$ oluşturmaktadır. Şekil~\ref{fig:denetim-modulasyon-blok} üzerindeki $v_{\mathrm{mod}}^{\ast}(t)$ ile Sin\textsuperscript{3} bloğunun $v_{\sin^3}(t)$ çıkışı gibi küçük harfli işaretler, zaman alanındaki fiziksel/normalize komutlardır. Eşitlik~\eqref{eq:modulasyon-toplami} içindeki büyük harfler ise aynı işaretlerin Laplace dönüşümleridir. Sin\textsuperscript{3} teriminin başarısı yalnızca düşük güç THD değeriyle değil, akım izleme hatası, kararlılık marjı, denetim eforu, doyum süresi ve anma akımındaki etkisiyle birlikte değerlendirilmektedir.

\subsection{Denetim ve kapı mantığının benzetim modelinden okunması}
\label{sec:denetim-modelden}

Denetim ve kapı mantığının tezdeki tanımı, deney senaryolarını yeniden üretmek için hazırlanan Simulink modellerine dayanmaktadır. Şekil~\ref{fig:denetim-modulasyon-blok} içindeki her blok, modelin blok yapısı ve sinyal bağlantıları okunarak doğrulanmıştır. Modelden okunan zincirin adımları aşağıdaki gibidir. Birim genlikli referans sinüsü $\sin(\theta+\varphi)$, \SI{1}{\micro\second} adımla örneklenmektedir. Sinüs, senaryoya göre $K_{1x}$ ile, S3'te ayrıca $K_{33}\sin^3$ dalıyla ölçeklenmekte ve görev oranı komutu $u$ oluşturulmaktadır. Yarım periyotlar, $u\geq0$ koşuluna bakan iki karşılaştırma anahtarıyla ayrılmıştır. Pozitif yarıda $d=u$, negatif yarıda $d=-u$ alınmaktadır. $d$, \SI{10}{\kilo\hertz} ve 0--1 aralıklı testere dişi taşıyıcıyla karşılaştırılmaktadır. Karşılaştırmada $V_{\mathrm{dc}}$ normalizasyonu yoktur. Dört kapı komutunun her biri, yalnız yükselen kenarı $t_d=\SI{1}{\micro\second}$ geciktiren bir gecikme bloğundan geçmektedir. Klasik anahtarlamada nötr anahtar komutları, dış anahtar PWM çıkışlarının mantıksal tümleyenidir. Susturmalı anahtarlamada ise nötr anahtarlar yarım periyot boyunca sabit 1 veya 0 değerini almaktadır. Kapı mantığının senaryo başına doğruluk tablosu Alt Bölüm~\ref{sec:model-kapi-mantigi} içinde verilmiştir.

Donanımın PWM zamanlayıcı yazmaçlarından periyot ve ölü zaman değerleri okunmuş ve modele aktarılmıştır (Tablo~\ref{tab:sistem-parametreleri}). S4 akım denetimi analog eşdeğer olarak modellenmiştir. Susturma ve Sin\textsuperscript{3} terimi, modellerde senaryo boyunca sürekli etkindir.

\section{Elektriksel ve Denetim Parametreleri}

Benzetim modellerinin ve deney düzeneğinin elektriksel, yarı iletken, modülasyon, denetim ve değerlendirme parametreleri Tablo~\ref{tab:sistem-parametreleri} altında tek yerde toplanmıştır. Tezin diğer bölümleri parametre değerini yinelememekte, tabloya başvurmaktadır. Kaynak sütunundaki kısaltmalar tablo notunda açıklanmıştır. Tür sütunu her değerin dayanağını altı sınıftan birine bağlamaktadır. \emph{Ölçülmüş} değer, deney düzeneğinin güç analizörü ya da ölçüm ekranından okunan çalışma noktasıdır. \emph{Veri sayfası} değeri, üretici veri sayfasından okunan ya da veri sayfası eğrisine oturtulan cihaz büyüklüğüdür. \emph{Tasarım} değeri, parametre dosyasında ya da donanım tasarım belgelerinde tanımlanan büyüklüktür. \emph{Kalibre} değer, benzetim modelini deney özetinin aktif ve reaktif gücüne oturtmak için ayarlanmıştır. \emph{Benzetim seçimi} yalnız benzetime ait modelleme ya da denetim kararını göstermektedir. \emph{Tez tanımı} değeri, tez tanımından türetilen ve benzetim modelinin doğrulanmasında belirlenen büyüklüktür (Tablo~\ref{tab:model-duzeltme-ozeti}).

\begingroup
\setlength{\tabcolsep}{4pt}
\begin{longtable}{>{\footnotesize\raggedright\arraybackslash}p{0.21\textwidth} >{\footnotesize\raggedright\arraybackslash}p{0.225\textwidth} >{\footnotesize\raggedright\arraybackslash}p{0.08\textwidth} >{\footnotesize\raggedright\arraybackslash}p{0.25\textwidth} >{\footnotesize\raggedright\arraybackslash}p{0.135\textwidth}}
\caption{Elektriksel, yarı iletken, modülasyon, denetim ve değerlendirme parametrelerinin birleşik tablosu.}\label{tab:sistem-parametreleri}\\
\toprule
Parametre & Değer & Birim & Kaynak & Tür \\
\midrule
\endfirsthead
\caption[]{Elektriksel, yarı iletken, modülasyon, denetim ve değerlendirme parametrelerinin birleşik tablosu (devamı).}\\
\toprule
Parametre & Değer & Birim & Kaynak & Tür \\
\midrule
\endhead
\bottomrule
\endfoot
%
\multicolumn{5}{@{}l}{\small\textbf{Sistem tanımı ve çalışma noktaları}} \\
Prototip faz yapısı & Üç faz, dört telli; nötr hattı DA bara orta noktasına bağlı & -- & ŞM & Tasarım \\
Benzetim ve deney kapsamı & Tek faz kolu; nötr dönüşü DA bara orta noktasına & -- & BM; DÖ tek faz büyüklükleri & Benzetim seçimi \\
Tasarım güç aralığı & 20--25 & \si{\kilo\watt} & TB & Tasarım (anma akımı bölgesi S4 benzetimiyle incelenmiştir) \\
Prototip anma görünür gücü & 20 & \si{\kilo\volt\ampere} & TB, akım sensörü boyutlandırma hesabı & Tasarım \\
Faz akımı tasarım değeri; üst sınırı & Yaklaşık 30; yaklaşık 36 & \si{\ampere} & \SI{20}{\kilo\volt\ampere} için faz başına güç hesabı; \%20 tasarım payı & Tasarım \\
Şebeke anma gerilimi & 230 faz--nötr; 400 faz--faz & \si{\volt} & TB & Tasarım \\
Şebeke frekansı $f_g$ & \num{50} & \si{\hertz} & PD & Tasarım \\
Deney şebeke gerilimi $V_g$ (RMS, güvenlik amacıyla varyakla sınırlandırılmış) & S1 \num{119.80}; S2 \num{121.00}; S3 \num{124.30} & \si{\volt} & DÖ; PD & Ölçülmüş \\
S4 şebeke gerilimi $V_g$ (RMS) & \num{120.0} & \si{\volt} & PD & Benzetim seçimi \\
DA bara gerilimi $V_{\mathrm{dc}}$ & S1 \num{431.5}; S2 \num{395.0}; S3 \num{427.5} & \si{\volt} & DÖ; PD & Ölçülmüş \\
S4 DA bara gerilimi $V_{\mathrm{dc}}$ & \num{427.5} & \si{\volt} & PD (S3 noktası) & Benzetim seçimi \\
DA bara tanımı & $V_{\mathrm{dc}}$: DA$+$ ile DA$-$ arasındaki toplam gerilim; yarım bara $V_{\mathrm{dc}}/2$ & -- & BM: tek DA kaynağı, iki seri bara kapasitörü, orta nokta nötre bağlı & Benzetim seçimi \\
Deney faz akımı $I_g$ (RMS) & S1 \num{2.26}; S2 \num{2.33}; S3 \num{2.31} & \si{\ampere} & DÖ; PD & Ölçülmüş \\
S4 akım referansı $I_{\mathrm{ref}}$ (RMS) & \num{2.3}; \num{10}; \num{25} & \si{\ampere} & PD (\SI{25}{\ampere}); düşük, orta ve anma akımı taraması & Benzetim seçimi \\
\midrule
\multicolumn{5}{@{}l}{\small\textbf{Ölçüm ve sayısal denetim donanımı}} \\
DA bara ölçüm aralığı & \num{470} & \si{\volt} & TB, AMC1200 giriş bölücüsü & Tasarım \\
Gerilim ölçümü & AMC1200, 12 bit ADC; ofset \SI{1.3}{\volt}, ADC referansı \SI{2.6}{\volt} & -- & TB, ölçüm kartı & Tasarım \\
Akım sensörü & HLSR 20-P/SP33; \SI{23}{\milli\volt\per\ampere}; ölçülebilir tepe \SI{56.5}{\ampere} & -- & VS; TB, ölçüm kartı & Veri sayfası \\
Aşırı akım eşiği & Yaklaşık 30,4 & \si{\ampere} & TB, pencere karşılaştırıcı eşikleri ve sensör kazancından hesap & Tasarım \\
Sayısal denetleyici & PIC32MK ailesi, 32 bit & -- & TB, kontrol kartı & Tasarım \\
\midrule
\multicolumn{5}{@{}l}{\small\textbf{LCL filtre, EMC süzgeci ve DA bara}} \\
Evirici tarafı endüktansı $L_1$ & \num{600} & \si{\micro\henry} & PD; ŞM & Tasarım \\
$L_1$ seri direnci & \num{0.03} & \si{\ohm} & PD & Tasarım \\
$L_1$ çekirdek kaybı paralel direnci $R_p$ & S1--S3 \num{600}; S4 \num{1500} & \si{\ohm} & PD & Tasarım \\
Filtre kapasitörü $C_f$ & \num{8} & \si{\micro\farad} & PD; ŞM & Tasarım \\
$C_f$ seri direnci & \num{0.01} & \si{\ohm} & PD & Tasarım \\
Şebeke tarafı endüktansı $L_2$ (şebeke ve varyak eşdeğeri) & \num{1000} & \si{\micro\henry} & PD & Tasarım (Alt Bölüm~\ref{sec:l2-tanimi}) \\
$L_2$ seri direnci & \num{0.01} & \si{\ohm} & PD & Tasarım \\
Ortak kip bobinleri (iki kademe) & Sargı \SI{400}{\micro\henry}; karşılıklı \SI{398}{\micro\henry} ($k=\num{0.995}$); sargı direnci \SI{5}{\milli\ohm}; çekirdek kaybı \SI{100}{\kilo\ohm}; diferansiyel eşdeğer kademe başına \SI{4}{\micro\henry} & -- & PD; ŞM ilk kademe \SI{460}{\micro\henry} & Tasarım \\
X kapasitörü ağı (tek faz eşdeğeri) & \SI{0.5}{\micro\farad}, \SI{0.2}{\ohm} (iki seri \SI{1}{\micro\farad}/\SI{0.1}{\ohm} kol) & -- & PD & Tez tanımı \\
DA bara kapasitörleri $C_{\mathrm{dc}+}=C_{\mathrm{dc}-}$ & \num{470} & \si{\micro\farad} & PD; ŞM \SI{470}{\micro\farad}/\SI{450}{\volt} & Tasarım (S1--S3, deney donanımı) \\
DA bara kapasitörleri, S4 orta ve anma akımı benzetimleri & \num{5600} & \si{\micro\farad} & Orta nokta salınımını \SI{25}{\ampere} akımda $V_{\mathrm{dc}}/2$ değerinin \%5'i altında tutan boyut & Benzetim seçimi \\
DA bara kapasitörü eşdeğer seri direnci & \num{0.10} & \si{\ohm} & PD & Tasarım \\
Bara dengeleme dirençleri & \num{220} & \si{\kilo\ohm} & PD & Tasarım \\
DA kaynak bağlantısı ve çıkış rölesi kontak direnci & \num{1} & \si{\milli\ohm} & PD & Tasarım \\
\midrule
\multicolumn{5}{@{}l}{\small\textbf{Yarı iletkenler ve kapı sürücüsü}} \\
Dış anahtarlar $S_1$--$S_2$ & SKM100GB12T4 & -- & ŞM; VS & Veri sayfası \\
$S_1$--$S_2$ iletim karakteristiği & $V_{CE}=\SI{0.9}{\volt}+\SI{11.5}{\milli\ohm}\cdot I_C$ (\SI{25}{\celsius}, $V_{GE}=\SI{15}{\volt}$) & -- & VS ($V_{CE0}$ ve $r_{CE}$ en büyük değerleri) & Veri sayfası \\
$S_1$--$S_2$ anahtarlama enerjisi (1--\SI{50}{\ampere}) & $E_{\mathrm{on}}$ 0,02--1,00; $E_{\mathrm{off}}$ 0,03--1,50 & \si{\milli\joule} & PD; VS eğilimine dayalı düşük akım ara değerlemesi (VS noktası \SI{600}{\volt}, \SI{100}{\ampere}, \SI{150}{\celsius}) & Tasarım \\
$S_1$--$S_2$ gecikme ve geçiş süreleri $t_{d(\mathrm{on})}$/$t_{d(\mathrm{off})}$/$t_r$/$t_f$ & 200/600/60/110 & \si{\nano\second} & PD (2--\SI{3}{\ampere} bölgesi) & Tasarım \\
Nötr anahtarlar $S_3$--$S_4$ & IXYN100N65C3H1 & -- & ŞM; VS & Veri sayfası \\
$S_3$--$S_4$ iletim karakteristiği & $V_{CE}$: \SI{0.80}{\volt} (\SI{1}{\ampere}) ile \SI{1.68}{\volt} (\SI{50}{\ampere}) arası tablo & -- & VS çıkış eğrisi (\SI{25}{\celsius}, $V_{GE}=\SI{15}{\volt}$) & Veri sayfası \\
$S_3$--$S_4$ anahtarlama enerjisi (1--\SI{50}{\ampere}) & $E_{\mathrm{on}}$ 0,05--1,30; $E_{\mathrm{off}}$ 0,07--0,83 & \si{\milli\joule} & PD; \SI{50}{\ampere} noktası VS, düşük akım bölgesi ara değerleme & Tasarım \\
$S_3$--$S_4$ gecikme ve geçiş süreleri $t_{d(\mathrm{on})}$/$t_{d(\mathrm{off})}$/$t_r$/$t_f$ & 200/50/100/1000 & \si{\nano\second} & PD & Tasarım \\
Dış anahtar ters paralel diyotları $DS_1$--$DS_2$ & $V_f=\SI{1.50}{\volt}$; $r_F=\SI{10.2}{\milli\ohm}$; $Q_{rr}=\SI{15.7}{\micro\coulomb}$; $I_{RM}=\SI{54}{\ampere}$ & -- & VS (en büyük değerler; \SI{100}{\ampere} toparlanma testi) & Veri sayfası \\
Nötr anahtar ters paralel diyotları $DS_3$--$DS_4$ & $V_f=\SI{1.70}{\volt}$; $r_F=\SI{6}{\milli\ohm}$; $Q_{rr}=\SI{9.5}{\micro\coulomb}$; $t_{rr}=\SI{100}{\nano\second}$ & -- & VS, iç diyot (\SI{100}{\ampere}; $Q_{rr}$ eğriden) & Veri sayfası \\
PFC yükseltici diyotları $D_1$--$D_2$ & DSEI2x61-12B; $V_f=\SI{1.58}{\volt}$; $r_F=\SI{8.5}{\milli\ohm}$; $t_{rr}=\SI{190}{\nano\second}$; $Q_{rr}=\SI{3.0}{\micro\coulomb}$; $C_j=\SI{31}{\pico\farad}$ & -- & VS ($t_{rr}$ için \SI{25}{\celsius} değeri; $Q_{rr}$ eğriden) & Veri sayfası \\
Kapı sürücüsü & $+15/-5$~\si{\volt}; toplam kapı direnci \SI{15}{\ohm}; çıkış yükselme ve düşme süresi \SI{20}{\nano\second}; ek açma ve kapama gecikmesi 0; IGBT kapı eşiği $V_{GE(th)}=\SI{6}{\volt}$ & -- & PD & Tasarım \\
Jonksiyon sıcaklığı & \num{25}, sabit & \si{\celsius} & PD & Benzetim seçimi \\
\midrule
\multicolumn{5}{@{}l}{\small\textbf{Modülasyon ve PWM}} \\
Anahtarlama frekansı $f_s$ & \num{10} & \si{\kilo\hertz} & PD; PWM zamanlayıcı periyodu (PTPER $=12000$) & Tasarım \\
Ölü zaman $t_d$ & \num{1} & \si{\micro\second} & PD (DTR $=$ ALTDTR $=120$); modelde her kapı komutunun yükselen kenarında gecikme & Tasarım \\
Taşıyıcı ve görev oranı & Testere dişi 0--1; yarım periyot ayrımlı tek kutuplu görev oranı $d=|u|$; $V_{\mathrm{dc}}$ normalizasyonu yok & -- & BM, PWM alt sistemi & Benzetim seçimi \\
Modülasyon referansı örnekleme adımı & \num{1} & \si{\micro\second} & PD & Benzetim seçimi \\
Temel modülasyon genliği $K_1$ & S1 \num{0.7865}; S2 \num{0.76796}; S3 \num{0.44343} & -- & KL & Kalibre \\
Sin\textsuperscript{3} kazancı $K_3$ & S3 \num{0.38431}; S1--S2 \num{0} & -- & KL; $K_3/K_1$ oranı sabit & Kalibre \\
Referans fazı $\varphi$ & S1 \num{-0.3661}; S2 \num{0.1121}; S3 \num{0.6575} & \si{\degree} & KL & Kalibre \\
S4 ileri besleme $K_1$, $K_3$, $\varphi$ & S3 kalibre değerleri & -- & KL & Kalibre \\
Sin\textsuperscript{3} teriminin biçimi & $K_{3x}\sin^3\theta$ (birim genlikli referansın küpü, pozitif işaretli) & -- & BM, modülasyon toplamı & Tasarım \\
Susturma ve Sin\textsuperscript{3} etkinleşme ölçütü & Senaryo boyunca sürekli etkin & -- & BM & Benzetim seçimi \\
\midrule
\multicolumn{5}{@{}l}{\small\textbf{S4 akım denetimi}} \\
Geri besleme büyüklüğü ve ölçüm süzgeci & $i_{L1}$; tek kutuplu alçak geçiren süzgeç, köşe \SI{5}{\kilo\hertz} ($f_s/2$) & -- & DT & Benzetim seçimi \\
Akım referansı & $\sqrt{2}\,I_{\mathrm{ref}}\sin\theta_g$; şebeke fazına kilitli, ideal faz bilgisi & -- & DT & Benzetim seçimi \\
Oransal kazanç $K_p$ & \num{0.03} & \si{\per\ampere} & PD & Tasarım \\
İntegral kazancı $K_i$ & \num{11.31} ($2\pi\cdot\SI{60}{\hertz}\cdot K_p$) & \si{\per\ampere\per\second} & DT; PI sıfırı \SI{60}{\hertz}, akım çevrimi kazanç geçiş frekansının (yaklaşık \SI{586}{\hertz}) onda biri & Benzetim seçimi \\
Geri hesaplamalı doyum önleme kazancı $K_{aw}$ & \num{377} ($2\pi\cdot60$) & \si{\per\second} & DT & Benzetim seçimi \\
PIR rezonans terimi & $K_r=\SI{0.3}{\per\ampere}$ ($10K_p$); $\omega_b=2\pi\cdot2$~\si{\radian\per\second}; $\omega_0=2\pi\cdot50$~\si{\radian\per\second} & -- & DT & Benzetim seçimi \\
Modülasyon ve integratör sınırı & $|u|\le1$ & -- & DT & Benzetim seçimi \\
Denetim adımı ve hesaplama gecikmesi & \SI{1}{\micro\second}; analog eşdeğer denetim & -- & DT & Benzetim seçimi \\
\midrule
\multicolumn{5}{@{}l}{\small\textbf{Değerlendirme}} \\
Değerlendirme penceresi ve THD tanımı & Son tam şebeke periyodu (\SI{20}{\milli\second}); THD, temel bileşene göre bir periyotluk kayan pencereyle & -- & BM, ortalama, RMS ve THD blokları (\SI{50}{\hertz}) & Benzetim seçimi \\
Güç ölçüm noktası ve işaret kuralı & Evirici çıkış uçları (EMC süzgeci sonrası, $L_2$ öncesi); evirici$\to$şebeke yönü pozitif & -- & BM, ölçüm zinciri & Benzetim seçimi \\
\midrule
\multicolumn{5}{@{}p{0.95\textwidth}@{}}{\footnotesize Kaynak kodları: PD parametre dosyası, VS üretici veri sayfası, ŞM tam topoloji şeması, DÖ deney özet tablosu. TB donanım ve ölçüm kartı tasarım belgelerini, BM benzetim modelini, DT denetleyici tasarımını, KL kalibrasyonu (Bölüm~\ref{chap:deney}) göstermektedir. Modelde laboratuvar kaynağının akım çekememesini temsil eden bir ters akım diyotu ($V_f=\SI{0.05}{\volt}$, $R_{\mathrm{on}}=\SI{1}{\ohm}$) yer almaktadır. SKM modülündeki \SI{4}{\micro\farad} bastırma kapasitörü ile \SI{26}{\nano\farad} Y kapasitörü modele katılmamıştır. Benzetimler ode15s değişken adımlı çözücüyle, \SI{1}{\micro\second} en büyük adımla ve \num{1e-6} bağıl, mutlak ve Simscape artık toleransıyla yürütülmüştür. Benzetim süresi S1--S3'te \SI{0.1}{\second}, S4'te P ve PI için \SI{0.06}{\second}, PIR için \SI{0.2}{\second}'dir.} \\
\end{longtable}
\endgroup

Tablodaki yarı iletkenlerden $D_1$--$D_2$ diyotları faz kolunda faz--DA rayı konumundadır ve GT-tipi mimarinin şarj kipinde Vienna PFC yükseltici diyotu olarak çalışmaktadır. Diyotlar kesim durumunda toplam $V_{\mathrm{dc}}$ gerilimini taşıdığından \SI{1200}{\volt} sınıfında harici diyot gerekmektedir. Nötr kol diyotları ise en çok $V_{\mathrm{dc}}/2$ gerilimini kestiğinden \SI{650}{\volt} sınıfındaki nötr IGBT'nin iç diyotu yeterlidir. Faz gerilimi DA raylarını aşmadığından kararlı çalışmada $D_1$--$D_2$ iletime girmemektedir. Sıfır durumunun serbest dolaşım akımının yolu $DS_3$--$DS_4$ diyotlarıdır. Bütün anahtarlamalı benzetimler, iletim karakteristiğini ve anahtarlama enerjisini akıma bağlı tablolarla veren veri sayfası tabanlı cihaz modelleriyle yapılmıştır. Benzetimler \SI{25}{\celsius} sabit jonksiyon sıcaklığında yürütülmüştür.

\subsection{Şebeke tarafı endüktansı $L_2$}
\label{sec:l2-tanimi}

LCL filtrenin şebeke tarafı endüktansı ayrı bir filtre bobini değildir. Deney düzeneğinde evirici çıkışı şebekeye, gerilimi güvenlik amacıyla sınırlandıran varyak üzerinden bağlanmıştır. $L_2=\SI{1000}{\micro\henry}$, şebeke empedansı ile varyak kaçak endüktansının toplam eşdeğeri olarak parametre dosyasında tanımlanmıştır ve tasarım türünde bir değerdir (Tablo~\ref{tab:sistem-parametreleri}). Tam topoloji şemasındaki \SI{600}{\micro\henry} etiketleri ise üç fazın $L_1$ bobinleridir. Donanımdaki fiziksel süzgeç, EMC süzgeciyle birlikte $L_1$--$C_f$ yapısındadır.

$L_2$ değerinin modelde iki etkisi vardır. İdeal LCL rezonansı, Eşitlik~\eqref{eq:lcl-rezonans} ile $L_1=\SI{600}{\micro\henry}$, $C_f=\SI{8}{\micro\farad}$ ve $L_2=\SI{1000}{\micro\henry}$ için yaklaşık \SI{2.9}{\kilo\hertz}'dir. EMC süzgecini ve kayıpları içeren fiziksel modelde rezonans \SI{2.8}{\kilo\hertz} çevresindedir (Alt Bölüm~\ref{sec:kararlilik-bode}). Rezonans frekansı anahtarlama frekansının altıda birinden büyük olduğundan $L_2$, akım denetiminin gecikme duyarlılığını belirleyen bir tasarım parametresidir. Şebeke endüktansının \num{0}--\SI{5}{\milli\henry} aralığındaki etkisi de kararlılık analizinde taranmıştır. Açık çevrim deney senaryolarında ise $L_2$, C noktasındaki gerilim ile şebeke iç gerilimi arasındaki küçük temel frekans farkını belirlemektedir. Kalibre edilen aktif ve reaktif güç değerleri, şebeke tarafı endüktansından yalnızca ikinci mertebede etkilenmektedir.

\section{Benzetim ve Deneysel Doğrulama Metodolojisi}
\label{sec:dogrulama-metodolojisi}

Önerilen anahtarlama ve kompanzasyon yöntemleri aynı model, çalışma noktası ve ölçüm tanımları altında karşılaştırılmıştır. Temel karşılaştırma kolları klasik tamamlayıcı anahtarlama (S1), nötr anahtar susturma yöntemi (S2) ve Sin\textsuperscript{3} kompanzasyonlu nötr anahtar susturma yöntemidir (S3). Deneylerde her kol için ayrı bir DA bara ve şebeke gerilimi ölçülmüştür. Her senaryonun benzetimi, ölçülen çalışma noktasında yapılmıştır. Dördüncü kol (S4), S3'e $L_1$ akım geri beslemesi eklenerek kurulan kapalı çevrim benzetimidir.

Tablo~\ref{tab:dogrulama-matrisi}, üç akım bölgesinde (düşük, orta ve anma) karşılaştırılan yöntemleri, her bölgenin kanıt türünü ve sonucun raporlandığı bölümü birleştirmektedir. Deneyler düşük güç noktasında yürütülmüştür. Orta ve anma akımı bölgesi ile geçici rejim, kalibre benzetim modelinden türetilen S4 benzetimiyle incelenmiştir. Benzetim ve deney sonuçları ayrı raporlanmaktadır. Kalibre benzetim, deneyin kalibre edilmeyen büyüklüklerini hangi hatayla yeniden ürettiği üzerinden değerlendirilmektedir.

\begin{table}[H]
  \centering
  \caption{Karşılaştırmalı doğrulama matrisi: çalışma bölgesi, yöntem, kanıt türü ve raporlandığı bölüm.}
  \label{tab:dogrulama-matrisi}
  \small
  \begin{tabularx}{\textwidth}{p{0.19\textwidth} p{0.22\textwidth} X p{0.17\textwidth}}
    \toprule
    Çalışma bölgesi & Yöntemler & Kanıt türü ve ölçütler & Bölüm \\
    \midrule
    Düşük güç (\SI{2.3}{\ampere}) & S1, S2, S3 & Deney özet tablosu ve ekranları ile kalibre benzetim. Verim, THD, $P_{\mathrm{dc}}$, $P_{\mathrm{ac}}$, $|Q_{g,1}|$, cos$\varphi_1$ & Bölüm~\ref{chap:anahtarlama}, Bölüm~\ref{chap:deney} \\
    Düşük güç, kapalı çevrim & S4: P, PI & Benzetim. Akım izleme, THD, doyum & Alt Bölüm~\ref{sec:s4-kapali-cevrim} \\
    Orta akım (\SI{10}{\ampere}) & S4: P, PI, PIR & Benzetim. Akım izleme, üçüncü harmonik, THD & Alt Bölüm~\ref{sec:s4-kapali-cevrim} \\
    Anma akımı (\SI{25}{\ampere}) & S4: P, PI, PIR & Benzetim. Akım izleme, THD, doyum, verim & Alt Bölüm~\ref{sec:s4-kapali-cevrim} \\
    Geçici rejim & S4: PI, \num{10}--\SI{25}{\ampere} basamağı & Benzetim. Aşım, yerleşme, geçişte doyum & Alt Bölüm~\ref{sec:s4-kapali-cevrim} \\
    Kararlılık & S4: P, PI, PIR; $K_3$ açık ve kapalı & Ortalama model ve anahtarlamalı model doğrusallaştırması. $L$, $S$, $T$, marjlar, kutuplar & Alt Bölüm~\ref{sec:kararlilik-bode} \\
    \bottomrule
  \end{tabularx}
\end{table}

Karşılaştırmaların aynı tanımlarla tekrarlanabilmesi için pencere başlangıcı $t_0$ ve kararlı durum analiz süresi $T_w$ olarak tanımlansın. DA bara gerilimi ve DA kaynak akımı sırasıyla $v_{\mathrm{dc}}$ ve $i_{\mathrm{dc}}$ ile gösterilsin. Evirici çıkış uçlarındaki (C noktası) gerilim ve şebeke akımı da sırasıyla $v_g$ ve $i_g$ olsun. Tek faz kolu için ortalama DA ve AA güçleri
\begin{equation}
  \begin{aligned}
    P_{\mathrm{dc}} &= \frac{1}{T_w}\int_{t_0}^{t_0+T_w}v_{\mathrm{dc}}(t)\,i_{\mathrm{dc}}(t)\,\mathrm{d}t, \\
    P_{\mathrm{ac}} &= \frac{1}{T_w}\int_{t_0}^{t_0+T_w}v_g(t)\,i_g(t)\,\mathrm{d}t
  \end{aligned}
  \label{eq:guc-ortalama-tanimlari}
\end{equation}
olarak hesaplanmaktadır. Eşitlik~\eqref{eq:guc-ortalama-tanimlari} içinde $i_{\mathrm{dc}}$, DA kaynağından bara pozitif ucuna akan akımdır. $i_g$ ise üretici kuralıyla, evirici fazından şebekeye doğru pozitif alınmaktadır. Deşarj yönünde $P_{\mathrm{dc}}>0$, şebekeye aktarılan güç için de $P_{\mathrm{ac}}>0$ olmaktadır. Benzetimde $T_w$ bir şebeke periyodudur (\SI{20}{\milli\second}). S1--S3 için değerler \SI{0.1}{\second} anındaki, yani 0,08--\SI{0.10}{\second} penceresinin ortalamasıdır. S4 için son tam şebeke periyodu kullanılmıştır. Deşarj yönündeki yüzde verim
\begin{equation}
  \eta=100\frac{P_{\mathrm{ac}}}{P_{\mathrm{dc}}}
  \label{eq:verim-tanimi}
\end{equation}
biçiminde raporlanmaktadır. Eşitlik~\eqref{eq:verim-tanimi}, yalnızca aynı zaman penceresinden elde edilen giriş ve çıkış güçlerine uygulanmaktadır. Reaktif güç, C noktasındaki gerilim ve akımın temel bileşen fazörlerinden $Q_{g,1}=\operatorname{Im}\{V_{g,1}I_{g,1}^{\ast}\}$ olarak hesaplanmaktadır. $Q_{g,1}>0$, akımın gerilimin gerisinde kaldığı ve evirici ile filtrenin şebekeye reaktif güç verdiği durumdur. Temel bileşen güç faktörü ($\cos\varphi_1$) ile toplam güç faktörü $\mathrm{PF}=P_{\mathrm{ac}}/(V_gI_g)$ ayrı büyüklüklerdir. İkisi arasındaki farkın kaynağı akım bozulmasıdır.

Akımın $h$'inci harmonik RMS bileşeni $I_h$, temel bileşen RMS değeri $I_1$ ve hesaba katılan en büyük harmonik mertebesi $H$ ile gösterilsin. Akım toplam harmonik bozunumu
\begin{equation}
  \mathrm{THD}_i=100\frac{\sqrt{\sum_{h=2}^{H}I_h^2}}{I_1}
  \label{eq:thd-rmse-tanimlari}
\end{equation}
ile belirlenmektedir \citep{Teodorescu2011,IEC61000}. S1--S3 benzetiminde THD, modelin \SI{50}{\hertz} temel frekanslı ve bir periyot uzunluklu THD bloğundan alınmıştır. S4 benzetiminde ise aynı pencerenin $H=50$ mertebesine kadar Fourier analizi kullanılmıştır. Deney tarafında THD değeri güç analizörü özetinden okunmuştur.

Verim hesabında DA giriş gücü ile AA çıkış gücü eşzamanlı veya zaman farkı açıkça sınırlandırılmış kayıtlar üzerinden karşılaştırılmalıdır. Deney altyapısında DA ve AA analizörleri arasındaki en yüksek kayıt gecikmesi yaklaşık \SI{50}{\milli\second}, haberleşme çevrimi ise \SI{1}{\second} düzeyindedir. Haberleşme kayıtları kararlı durum verimi için uygundur. Geçici rejim, S4 benzetimiyle incelenmiştir.

Benzetim tarafındaki kanıt zinciri izlenebilirdir. Her sonuç, kendisini üreten betiğe ve benzetim çıktılarına bağlıdır. Aynı betik aynı sonuçları yeniden üretmektedir. Deney tarafındaki kanıt ise özet tablo ve ölçüm ekranlarından oluşmaktadır.

\section{Benzetim Modelinin Doğrulanması}
\label{sec:model-dogrulama-duzeltme}

Benzetim modelleri, deney senaryolarını yeniden üretmek amacıyla senaryo başına hazırlanmıştır. S1 ile S2 ortak bir model dosyasında iki bağımsız devre olarak, S3 ise ayrı bir dosyada yer almaktadır. Dört senaryo ortak bir parametre dosyasını okumaktadır. Modeller, tezde kullanılmadan önce kalite güvence adımı olarak üç denetimden geçirilmiştir. Birinci denetimde modellerin benzetim sonuçları deney özet tablosuyla karşılaştırılmıştır. İkinci denetimde güç devresinin bağlantı grafiği, bloklar arası çizgilerden yeniden kurulmuştur. Her IGBT ve diyot bloğunun devre rolü, parametre kümesi ve kapı etiketi tam topoloji şemasıyla ve senaryo tanımlarıyla karşılaştırılmıştır. Üçüncü denetimde ölçüm zinciri ve güç işaret kuralı ham dalga biçiminden doğrulanmıştır.

Denetimler sonucunda modelde uygulanan doğrulama adımları Tablo~\ref{tab:model-duzeltme-ozeti} içinde özetlenmiştir. Deney değerlerine uyum, Alt Bölüm~\ref{sec:kalibrasyon} içinde anlatılan iki serbest değişkenli kalibrasyonla sağlanmıştır. Doğrulamada belirlenen parametreler Tablo~\ref{tab:sistem-parametreleri} içinde ``Tez tanımı'' türüyle işaretlidir.

\begingroup
\setlength{\tabcolsep}{3.5pt}
\begin{longtable}{>{\footnotesize\raggedright\arraybackslash}p{0.235\textwidth} >{\footnotesize\raggedright\arraybackslash}p{0.21\textwidth} >{\footnotesize\raggedright\arraybackslash}p{0.235\textwidth} >{\footnotesize\raggedright\arraybackslash}p{0.235\textwidth}}
\caption{Benzetim modelinin doğrulama adımları: denetlenen özellik, hedef durum, model uygulaması ve sonuç.}\label{tab:model-duzeltme-ozeti}\\
\toprule
Denetlenen özellik & Hedef durum & Model uygulaması & Sonuç \\
\midrule
\endfirsthead
\caption[]{Benzetim modelinin doğrulama adımları: denetlenen özellik, hedef durum, model uygulaması ve sonuç (devamı).}\\
\toprule
Denetlenen özellik & Hedef durum & Model uygulaması & Sonuç \\
\midrule
\endhead
\bottomrule
\endfoot
\textbf{S2 kapı yolu.} Dört kapı girişinin susturmalı PWM ile sürülmesi & Nötr IGBT'ler yarım periyot boyunca 1/0; genlik ve faz S2 parametrelerinden & S2 kapı işaretleri, S3'ün susturmalı PWM alt sisteminden $K_{33}=0$ ile üretilmektedir & $\eta=\%94{,}07$ (deney \%94,2); kayıp \SI{16.99}{\watt} (deney \SI{16.98}{\watt}) \\
\textbf{Diyot parametre eşlemesi.} Diyot parametre kümelerinin devre rolüyle eşlenmesi & $D_1$--$D_2$ için \SI{1200}{\volt} harici diyot, $DS_3$--$DS_4$ için nötr IGBT iç diyotu & Sekiz diyot bloğu, rolüne uygun parametre kümesini kullanmaktadır & Eşlemenin etkisi S3: $P_{\mathrm{ac}}$ $-\%0{,}15$, $\eta$ $-0{,}04$ puan; S1: $P_{\mathrm{ac}}$ $-\%0{,}9$, $\eta$ $+0{,}07$ puan \\
\textbf{X kapasitörü eşdeğeri.} İki seri X kapasitörü kolunun tek faz eşdeğeri & İki seri \SI{1}{\micro\farad} kolun eşdeğeri \SI{0.5}{\micro\farad}/\SI{0.2}{\ohm} & İki X kapasitörü bloğu eşdeğer değerleri kullanmaktadır & Eşdeğerin etkisi S3: $P_{\mathrm{ac}}$ $+\%0{,}31$, $\eta$ $+0{,}006$ puan \\
\end{longtable}
\endgroup

Diyot parametre eşlemesinin aktif güce etkisi S1'de S3'tekinden belirgin biçimde büyüktür (\%0,9'a karşı \%0,15). S1'de nötr kol tamamlayıcı anahtarlandığından $DS_3$--$DS_4$ her anahtarlama periyodunda ters toparlanmaya girmektedir. S3'te ise nötr IGBT'ler yarım periyot boyunca iletimde kaldığından toparlanma olayı seyrektir.

\section{Düşük Akımda Benzetim Modelleme}
\label{sec:dusuk-akim-modelleme}

Laboratuvar çalışma noktası yaklaşık \SI{2.3}{\ampere} RMS, yani yaklaşık \SI{3.3}{\ampere} tepe faz akımıdır ve faz akımı tasarım değerinin \%8'i düzeyindedir (Tablo~\ref{tab:sistem-parametreleri}). Cihaz modellerinin veri sayfası noktaları ise \num{50}--\SI{100}{\ampere} bölgesindedir. Düşük akım bölgesinde modelin doğruluğunu dört etki belirlemektedir: anahtarlama enerjisinin düşük akıma ara değerlemesi, diyot ters toparlanması, akım dalgalanmasının akım genliğine oranı ve susturmayla oluşan süreksiz iletim. Aşağıdaki paragraflar dört etkiyi sırasıyla ele almaktadır.

\paragraph{Anahtarlama enerjisinin ara değerlemesi.} Dış ve nötr IGBT'lerin açma ve kapama enerjileri \num{1}--\SI{50}{\ampere} aralığında tablolarla tanımlanmıştır. \SI{50}{\ampere} noktası veri sayfasından, düşük akım bölgesi ise veri sayfası eğilimine dayalı ara değerlemeden gelmektedir (Tablo~\ref{tab:sistem-parametreleri}). Ara değerlemenin sonuca etkisi, kayıp dağılımı üzerinden sınırlanabilir. S3 modelinin kalibrasyon öncesi çalışma noktasında ($P_{\mathrm{ac}}\approx\SI{298}{\watt}$) blok bazında ayrıştırılan toplam kayıp \SI{18.94}{\watt}'tır. Kaybın \SI{11.56}{\watt}'ı $L_1$ çekirdek kaybı direncindedir. Dış anahtarların iletim kaybı anahtar başına yaklaşık \SI{1.37}{\watt}, anahtarlama kaybı \num{0.65}--\SI{0.66}{\watt}'tır. Nötr IGBT'lerin iletim kaybı anahtar başına \SI{0.41}{\watt}, $DS_3$--$DS_4$ diyotlarının kaybı ise diyot başına \SI{0.63}{\watt}'tır. Anahtarlama kaybı toplam kaybın yaklaşık \%7'si olduğundan ara değerlemedeki bir belirsizlik verimi ancak aynı oranda etkilemektedir. Düşük akım veriminde belirleyici parametre $L_1$ çekirdek kaybı direncidir. S4 benzetimleri \SI{1500}{\ohm} çekirdek kaybı direnciyle yürütüldüğünden S4 verimi, S4 akım taraması içinde değerlendirilmektedir.

\paragraph{Diyot ters toparlanması.} Sıfır durumunun serbest dolaşım akımı $DS_3$--$DS_4$ üzerinden akmaktadır. S3'te diyot başına ortalama akım \SI{0.36}{\ampere}, RMS akım \SI{1.17}{\ampere}'dir. Nötr IGBT'nin iç diyotunun ters toparlanma tepe akımı \SI{-3.4}{\ampere} düzeyindedir ve düşük akımda toparlanma enerjisine katkısı küçüktür. Klasik anahtarlamada (S1) ise nötr kol her periyotta tamamlayıcı anahtarlandığından diyot toparlanması her anahtarlama periyodunda tekrarlanmaktadır. Diyot parametre eşlemesinin aktif güce etkisinin S1'de S3'tekinden belirgin biçimde büyük olması (Tablo~\ref{tab:model-duzeltme-ozeti}), toparlanma parametrelerinin klasik anahtarlamada belirleyici olduğunu göstermektedir. Simscape diyot modelinin anahtarlama kenarlarında ürettiği ani akım tepeleri parazitik endüktanssız ideal devreye aittir. Diyot zorlanması, yukarıdaki ortalama ve RMS akım değerleriyle verilmiştir.

\paragraph{Akım dalgalanması ve akım genliği.} Eşitlik~\eqref{eq:buck-akim-dalgalanmasi}'na göre pozitif yarım periyotta $L_1$ akımının tepe--tepe dalgalanması $v_g=V_h/2$ anında en büyük değerine ulaşmaktadır: $\Delta i_{L,\max}=V_hT_s/(4L_1)$. $V_h=V_{\mathrm{dc}}/2\approx\SI{214}{\volt}$, $T_s=\SI{100}{\micro\second}$ ve $L_1=\SI{600}{\micro\henry}$ için yaklaşık \SI{9}{\ampere} tepe--tepe dalgalanma elde edilmektedir. Eşitlik~\eqref{eq:ters-akim-esigi} gereği $|\bar i_L|\lesssim\SI{4.5}{\ampere}$ olan her bölgede akım bir anahtarlama periyodu içinde sıfıra ulaşmaktadır. Laboratuvar çalışma noktasının tepe akımı yaklaşık \SI{3.3}{\ampere}'dir. Dolayısıyla klasik anahtarlamada akım, temel periyodun büyük bölümünde her anahtarlama periyodunda işaret değiştirmektedir. Nötr kol ters akımı, sıfır geçişleri çevresiyle sınırlı kalmayan, düşük güç çalışmasının genel bir durumudur. Aynı ölçüt S4'te orta akım davranışını da açıklamaktadır. \SI{25}{\ampere}'de süreksiz iletim yalnız sıfır geçişleri çevresinde oluşurken \SI{10}{\ampere}'de temel periyodun büyük bölümünü kapsamaktadır (Alt Bölüm~\ref{sec:s4-kapali-cevrim}).

\paragraph{Süreksiz iletim ve senkron olmayan alçaltıcı davranışı.} Nötr anahtar susturulduğunda, örneğin pozitif yarım periyotta $Q_3$ kesimde tutulduğunda, sıfır durumunda yalnız yük akımı yönündeki yol açıktır. $L_1$ akımı sıfıra indiğinde ters yönde devam edememektedir. Köprü akımı sıfır kalmakta, kutup gerilimi ise üç seviyeden birine bağlanmayıp filtre tarafının gerilimini izlemektedir. Sonuç olarak faz kolu, senkron alçaltıcı dönüştürücüden senkron olmayan, diyotlu alçaltıcı dönüştürücüye geçmektedir. Süreksiz iletim, köprüyü modülasyon komutuyla belirlenen sert bir gerilim kaynağı yerine yumuşak bir kaynak gibi davrandırmaktadır. Kutup geriliminin ara seviyede kaldığı zaman oranı kalibrasyon öncesi benzetimlerde S1 için \%0,5 (yalnız ölü zaman geçişleri), S2 için \%12,4, S3 için \%22,6'dır. Şekil~\ref{fig:s1-s2-kalibre-dalga}(c), S1 ve S2'nin kutup gerilimini aynı anahtarlama periyotlarında karşılaştırmaktadır. S1'de kutup gerilimi $+V_{\mathrm{dc}}/2$ ile 0 arasında anahtarlanmaktadır. S2'de ise sıfır durumunun ardından akım söndüğünde kutup gerilimi ara bir seviyeye sıçramaktadır.

Benzetimde $L_1$ doğrusal endüktans olarak modellenmiştir. Kalibre S1 ve S2 modellerinin THD farkı ile S2 şebeke akımı farkı (Alt Bölüm~\ref{sec:kalibrasyon}), ikinci mertebe düşük akım etkileriyle ilişkilendirilebilir. İkinci mertebe etkiler arasında IGBT çıkış kapasitesi ve sürücü yayılım gecikmeleri, en küçük darbe sınırı, parazitik endüktanslar ve şebeke gerilimindeki harmonikler yer almaktadır.

\section{Bölüm Sonucu}

Sistem mimarisi güç, ölçüm ve denetim yollarını ayıracak biçimde tanımlanmıştır. Prototipin üç fazlı yapısı ile benzetim ve deneyin tek faz kolu kapsamı da ayrılmıştır. T-tipi faz kolunun üç gerilim seviyesi ve dört temel akım yolu ortak gösterim altında toplanmış, benzetim modelinden okunan kapı mantığının Tablo~\ref{tab:t-tipi-durumlar} ile örtüştüğü gösterilmiştir. Düşük modülasyon bölgesinde senkron-buck eşdeğeri üzerinden ters akım eşiği açıklanmış, LCL filtre için durum denklemleri ve ideal aktarım fonksiyonu kurulmuştur.

Bütün devre, yarı iletken, modülasyon ve denetim parametreleri değer, birim, kaynak ve tür bilgisiyle Tablo~\ref{tab:sistem-parametreleri} içinde toplanmıştır. Verim, güç, reaktif güç ve THD tek faz ölçüm noktasına göre tanımlanmıştır. Benzetim modeli tez tanımına ve tam topoloji şemasına göre doğrulanmıştır. Düşük akımda süreksiz iletim, köprüyü yumuşak bir kaynak gibi davrandırmaktadır. Kurulan model, sonraki bölümde önerilen anahtarlama yönteminin analizine ve Bölüm~\ref{chap:sin3} içindeki kapalı çevrim ile Bode incelemesine temel oluşturmaktadır.
